\documentclass[12pt]{article}
\usepackage[ngerman]{babel} 
\usepackage{amsmath, amssymb, amsthm, mathtools}
\usepackage{tikz}
\usepackage{xcolor}
\usepackage{mathrsfs}
\usepackage[margin=1in]{geometry}
\usepackage{parskip}
\usepackage{enumitem}
\usepackage{titlesec}

% ==========================================
% E-INK GRADIENT PALETTE V2 (Optimiert für Kontrast)
% ==========================================
\definecolor{eBlack}{RGB}{0,0,0}               
\definecolor{eMediumBlue}{RGB}{0,0,205}        
\definecolor{eDarkRed}{RGB}{139,0,0}           
\definecolor{eDarkGreen}{RGB}{0,100,0}         
\definecolor{eSaddleBrown}{RGB}{139,69,19}     
\definecolor{eDarkCyan}{RGB}{0,139,139}        
\definecolor{eOlive}{RGB}{128,128,0}           
\definecolor{eMediumOrchid}{RGB}{186,85,211}   
\definecolor{ePaleVioletRed}{RGB}{219,112,147} 
\definecolor{eSpringGreen}{RGB}{0,255,127}     
\definecolor{eSandyBrown}{RGB}{244,164,96}     
\definecolor{eSilver}{RGB}{192,192,192}        
\definecolor{eGold}{RGB}{255,215,0}            
\definecolor{eKhaki}{RGB}{240,230,140}         
\definecolor{eMistyRose}{RGB}{255,228,225}     

% Grauwerte
\definecolor{eDarkGray}{RGB}{64,64,64}   
\definecolor{eGray}{RGB}{128,128,128}    
\definecolor{eLightGray}{RGB}{211,211,211} 

% --- GLOBALE MATHEMATIK-FARBE ---
\everymath{\color{eMediumBlue}}
\everydisplay{\color{eMediumBlue}}

% --- TITEL-FORMATIERUNG ---
\titleformat{\section}{\huge\sffamily\bfseries\color{eMediumBlue}}{\thesection.}{0.5em}{}
\titleformat{\subsection}{\Large\sffamily\bfseries\color{eDarkCyan}}{\thesubsection}{0.5em}{}

% --- LISTEN-FORMATIERUNG ---
\setlist[enumerate,1]{label=\textbf{\textcolor{eDarkRed}{\alph*)}}}
\setlist[enumerate,2]{label=\textbf{\textcolor{eDarkRed}{\roman*)}}}
\setlist[itemize,1]{label=\textcolor{eDarkRed}{$\bullet$}}
\setlist[itemize,2]{label=\textcolor{eDarkRed}{$\circ$}}

% --- STANDARD MAKROS ---
\newcommand{\R}{\mathbb{R}}
\newcommand{\N}{\mathbb{N}}
\newcommand{\eps}{\varepsilon}
\newcommand{\ball}{\mathcal{B}}
\newcommand{\set}[1]{\left\{ #1 \right\}}
\newcommand{\warning}[1]{\textcolor{eDarkRed}{#1}}
\newcommand{\qt}[1]{\textit{„#1“}}
\DeclareMathOperator{\grad}{grad}

% --- GEMINI STYLES ---
\newcommand{\gemininote}[1]{\textcolor{eSandyBrown}{\textbf{[Gemini-Notiz: #1]}}}
\newcommand{\geminiwarning}[1]{\textcolor{eDarkRed}{\textbf{[Gemini-Warnung: #1]}}}

% Mathematische Operatoren
\DeclarePairedDelimiter\abs{\lvert}{\rvert}
\DeclarePairedDelimiter\norm{\lVert}{\rVert}
\DeclareMathOperator{\Span}{span}

% --- THEOREM-UMGEBUNGEN ---
\newtheoremstyle{taBlue}
  {12pt}{12pt}{\normalfont}{}{\bfseries\color{eMediumBlue}}{.}{0.5em}{}
\newtheoremstyle{taGreen}
  {12pt}{12pt}{\normalfont}{}{\bfseries\color{eDarkGreen}}{.}{0.5em}{}
\newtheoremstyle{taBrown}
  {12pt}{12pt}{\normalfont}{}{\bfseries\color{eSaddleBrown}}{.}{0.5em}{}
\newtheoremstyle{taRed}
  {12pt}{12pt}{\normalfont}{}{\bfseries\color{eDarkRed}}{.}{0.5em}{}

\theoremstyle{taBlue}
\newtheorem{theorem}{Theorem}[section]
\newtheorem{proposition}[theorem]{Proposition}
\newtheorem{proposition*}{Proposition}
\newtheorem*{theorem*}{Theorem}

\theoremstyle{taGreen}
\newtheorem{definition}[theorem]{Definition}
\newtheorem*{definition*}{Definition}
\newtheorem{example}[theorem]{Example}
\newtheorem*{example*}{Example}
\newtheorem{answer}{Answer}[section]
\newtheorem*{answer*}{Answer}

\theoremstyle{taBrown}
\newtheorem{exercise}{Exercise}[section]
\newtheorem*{exercise*}{Exercise}
\newtheorem*{exercises}{Exercises}
\newtheorem{question}{Question}[section]
\newtheorem*{question*}{Question}
\newtheorem*{warmup}{Warm-up}

\theoremstyle{taRed}
\newtheorem*{remark}{Remark}

\newenvironment{solution}{\begin{proof}[\textbf{\textcolor{eDarkGreen}{Solution}}]}{\end{proof}}


% --- COUNTERS FOR DEFINITION NUMBERING TRICK ---
\newcounter{tempSection}
\newcounter{tempTheorem}


\begin{document}
\setcounter{section}{0} % Start at section 2

% Week 2 - Monday: Metric Spaces, Open and Closed Sets, Sequences and Convergence
%\section{Metric Spaces}

A \textbf{metric space} $(X, d)$ is any set $X$ equipped with a \textbf{metric}. 

Formally, a metric is a function:
\[ d: X \times X \to \mathbb{R}_{\ge 0} \]
such that for all $x, y, z \in X$, the following three properties hold:

\begin{enumerate}
    \item \textbf{Definiteness:} $d(x,y) \ge 0$, with equality if and only if $x=y$. 
    \\ \textcolor{eGray}{In words: \qt{the distance is positive, and a point has zero distance to itself.}} 
    
    \item \textbf{Symmetry:} $d(x,y) = d(y,x)$.
    \\ \textcolor{eGray}{In words: \qt{$x$ has the same distance to $y$ as $y$ to $x$.}} 
    
    \item \textbf{Triangle Inequality:} $d(x,z) \le d(x,y) + d(y,z)$.
    \\ \textcolor{eGray}{In words: \qt{detours make the way longer.}} 
\end{enumerate}

\begin{center}
\begin{tikzpicture}[scale=1.5, very thick]
    % Triangle Inequality visual
    \coordinate (X) at (0,0);
    \coordinate (Y) at (2,1.5);
    \coordinate (Z) at (4,0);

    \draw[eMediumBlue, ->] (X) -- node[above left] {$d(x,y)$} (Y);
    \draw[eMediumBlue, ->] (Y) -- node[above right] {$d(y,z)$} (Z);
    \draw[eDarkRed, dashed, ->] (X) -- node[below] {$d(x,z) \le d(x,y) + d(y,z)$} (Z);

    \fill[eSaddleBrown] (X) circle (2pt) node[left] {$x$};
    \fill[eSaddleBrown] (Y) circle (2pt) node[above] {$y$};
    \fill[eSaddleBrown] (Z) circle (2pt) node[right] {$z$};
\end{tikzpicture}
\end{center}

\hrulefill

\subsection*{Examples of Metric Spaces}

\begin{example*}[The Discrete Metric]
Let $X$ be any set. The discrete metric is given by:
\[
d(x,y) = 
\begin{cases} 
1 & \text{if } x \ne y \\ 
0 & \text{if } x = y 
\end{cases}
\]
\textcolor{eGray}{Note: This is super important for finding counterexamples!}
\end{example*}

\begin{example*}[Metrics in $\R^n$]
We have many choices for a metric on $\R^n$. Some common ones are:
\begin{itemize}
    \item \textbf{Manhattan metric ($d_1$):} $d_1(x,y) = \sum_{k=1}^n \abs{x_k - y_k}$
    \item \textbf{Standard/Euclidean metric ($d_2$):} $d_2(x,y) = \sqrt{\sum_{k=1}^n (x_k - y_k)^2}$
    \item \textcolor{eDarkGreen}{\qt{Supremum metric}} \textbf{($d_3$ or $d_\infty$):} $d_3(x,y) = \sup_{k=1,\dots,n} \abs{x_k - y_k}$
\end{itemize}
\textcolor{eGray}{Generalization:} $d_m(x,y) = \sqrt[m]{\sum_{k=1}^n \abs{x_k - y_k}^m}$ is a metric on $\R^n$ for any $m \ge 1$.

\begin{center}
\textbf{Visualizing the Unit Balls in $\R^2$:} \\
\vspace{0.3cm}
\begin{tikzpicture}[scale=2]
    % Axes
    \draw[->, thick] (-1.5,0) -- (1.5,0) node[right] {$x_1$};
    \draw[->, thick] (0,-1.5) -- (0,1.5) node[above] {$x_2$};
    
    % d3: Supremum (Square)
    \draw[eDarkGreen, very thick, dashed] (-1,-1) rectangle (1,1);
    \node[eDarkGreen] at (1.2, 1.2) {$d_3$ (Supremum)};
    
    % d2: Euclidean (Circle)
    \draw[eMediumBlue, very thick] (0,0) circle (1cm);
    \node[eMediumBlue] at (1.2, 0.5) {$d_2$ (Standard)};
    
    % d1: Manhattan (Diamond)
    \draw[eDarkRed, very thick, dotted] (1,0) -- (0,1) -- (-1,0) -- (0,-1) -- cycle;
    \node[eDarkRed] at (0.7, -0.7) {$d_1$ (Manhattan)};
\end{tikzpicture}
\end{center}
\end{example*}

\begin{example*}[Continuous Functions]
Let $X = C^0([0,1], \R)$ denote the set of all continuous functions mapping the interval $[0,1]$ to $\R$. We can equip this space with the \textcolor{eDarkGreen}{\qt{Supremum Metric}}, defined for any two functions $f,g \in X$ as:
\[ d(f,g) := \sup_{x \in [0,1]} \abs{f(x) - g(x)} \]
\textcolor{eGray}{Intuition: This metric finds the single point $x$ where the two functions are furthest apart, and uses that maximum vertical gap as the distance between the two functions.}
\end{example*}

\begin{example*}[The n-Sphere]
Let $X = S^2 := \set{ (x_1, x_2, x_3) \in \R^3 : x_1^2 + x_2^2 + x_3^2 = 1 }$ (the 2-dimensional sphere). 
The closest path between two points on the sphere is an arc of a great circle (a geodesic). This defines a metric, similar to how straight lines define a metric on $\R^n$.

\begin{center}
\begin{tikzpicture}[scale=2]
    % Sphere Outline
    \draw[thick] (0,0) circle (1cm);
    \draw[dashed] (1,0) arc (0:180:1cm and 0.3cm);
    \draw (1,0) arc (0:-180:1cm and 0.3cm);
    
    % Points
    \coordinate (X) at (-0.7, 0.4);
    \coordinate (Y) at (0.6, 0.6);
    
    \fill[eSaddleBrown] (X) circle (1.5pt) node[left] {$x$};
    \fill[eSaddleBrown] (Y) circle (1.5pt) node[right] {$y$};
    
    % Euclidean Distance (straight line cutting through sphere)
    \draw[eDarkRed, dashed, thick] (X) -- node[below=0.2cm, sloped] {\small Euclidean distance} (Y);
    
    % Great Circle Distance (arc on surface)
    \draw[eMediumBlue, very thick] (X) to[bend left=20] node[above, sloped] {\small distance on sphere} (Y);
\end{tikzpicture}
\end{center}
\end{example*}

\hrulefill
\vspace{0.5cm}

\begin{exercise*}
Show that the supremum metric from Example 3 is indeed a metric.
\end{exercise*}

\begin{solution}
Symmetry and definiteness are clear. For the triangle inequality, let $f, g, h \in C^0([0,1])$ and set:
\[ p := f - h, \quad q := h - g \]
Then for any $x \in [0,1]$, we know that:
\[ \abs{p(x)} \le \sup \abs{p} \quad \text{and} \quad \abs{q(x)} \le \sup \abs{q} \]
Adding these inequalities, for all $x \in [0,1]$ we get:
\begin{equation}
    \begin{aligned}
        \abs{f(x) - g(x)} &= \abs{p(x) + q(x)} \\
        &\le \abs{p(x)} + \abs{q(x)} \quad \text{\textcolor{eGray}{(standard triangle inequality in $\R$)}} \\
        &\le \sup \abs{p} + \sup \abs{q} \\
        &= \sup \abs{f - h} + \sup \abs{h - g}
    \end{aligned}
\end{equation}
Since this final bound holds for \textcolor{eSaddleBrown}{every} $x \in [0,1]$, it must also hold for the supremum over all $x$. Thus:
\[ d(f,g) \le d(f,h) + d(h,g) \]
which proves the triangle inequality!
\end{solution}

%
\section{Open and Closed Sets}

Let $(X, d)$ be a metric space, let $x \in X$, and let $r \in \R_{>0}$. 
First, we define the \textcolor{eSaddleBrown}{\qt{open ball}} of radius $r$:
\[ B_r(x) := \set{ y \in X : d(y,x) < r } \]
\textcolor{eGray}{Note: Pay attention to the strict inequality!}

\begin{definition*}
\begin{itemize}
    \item A subset $U \subseteq X$ is \textcolor{eDarkGreen}{\textbf{open}} iff for every $y \in U$, there exists an $\eps > 0$ such that $B_\eps(y) \subseteq U$.
    \item A subset $A \subseteq X$ is \textcolor{eDarkRed}{\textbf{closed}} iff its complement $X \setminus A$ is open.
\end{itemize}
\end{definition*}

\begin{center}
\begin{tikzpicture}[scale=1.2, very thick]
    % Visualizing an open set
    \draw[eMediumBlue, fill=eMediumBlue!10, dashed] (0,0) .. controls (1,2) and (3,1) .. (4,0) .. controls (3,-2) and (1,-1) .. (0,0);
    \node[eMediumBlue] at (3.5, 1) {\Large $U$};
    
    % The open ball inside
    \coordinate (Y) at (1.5, 0);
    \draw[eDarkGreen, fill=eDarkGreen!20, dashed] (Y) circle (0.6cm);
    \fill[eSaddleBrown] (Y) circle (2pt) node[below] {$y$};
    \draw[eDarkGreen, ->] (Y) -- ++(45:0.6cm) node[midway, above left] {$\eps$};
    \node[eDarkGreen] at (2.5, -0.5) {$B_\eps(y) \subseteq U$};
\end{tikzpicture}
\end{center}

\vspace{0.5cm}

\begin{exercise*}
Let $(X,d)$ be a metric space. Show that $X$ and $\emptyset$ are both open and closed.
\end{exercise*}

\begin{solution}
First, let's look at openness:
\begin{itemize}
    \item $X$ is open because for any $x \in X$ and any $r > 0$, the ball $B_r(x)$ is by definition a subset of $X$. Thus, $B_r(x) \subseteq X$ trivially holds.
    \item $\emptyset$ contains no points. Therefore, the condition \qt{for every $y \in \emptyset$\dots} is satisfied vacuously. Thus, the empty set is open.
\end{itemize}
Now for closedness:
Since $X$ and $\emptyset$ are each other's complements in $X$ (i.e., $X \setminus X = \emptyset$ and $X \setminus \emptyset = X$), and we just proved both are open, they must both be closed as well!
\end{solution}

\hrulefill

%
\section{Sequences and Convergence}

A sequence $(x_n)_{n=0}^\infty$ in a metric space $(X, d)$ is formally a function $x: \N \to X$. We say that $(x_n)$ \textbf{converges} to $x \in X$ if:
\[ \forall \eps > 0, \exists N \in \N \text{ s.t. } d(x_n, x) < \eps \quad \forall n \ge N \]

\subsection*{The Big Connection Proposition}
Topological properties and sequences go hand-in-hand:
\begin{enumerate}
    \item $U \subseteq X$ is \textbf{open} $\iff$ for every sequence in $X$ with a limit in $U$, the sequence eventually lies inside $U$.
    \item $A \subseteq X$ is \textbf{closed} $\iff$ every convergent sequence in $A$ has its limit in $A$.
\end{enumerate}

\subsection*{\warning{Watch out! A Common Misconception!}}
Let's look at the standard Euclidean metric space $\R^2$. 
We know $[0,1]^2$ is closed and $(0,1)^2$ is open. What about $[0,1] \times (0,1)$? It is neither open nor closed in $\R^2$.

\textcolor{eDarkRed}{But...} what if our \textit{entire universe} (our metric space $X$) is strictly just $[0,1]^2$? 
In the subspace $X = [0,1]^2$, the set $[0,1]^2$ is actually \emph{open}! 

\textcolor{eDarkCyan}{Why?} Because the full space $X$ is always open. If you take a point on the boundary, the open ball $B_\eps(x)$ only consists of points \textcolor{eSaddleBrown}{that actually exist in $X$}. It doesn't bleed out into $\R^2$ because $\R^2$ doesn't exist to this metric space!

\begin{center}
\begin{tikzpicture}[scale=2, very thick]
    % Misconception visual
    \draw[->] (-0.5, 0) -- (1.5, 0) node[right] {$x_1$};
    \draw[->] (0, -0.5) -- (0, 1.5) node[above] {$x_2$};
    
    % The space [0,1]^2
    \draw[eMediumBlue, fill=eMediumBlue!10] (0,0) rectangle (1,1);
    \node[eMediumBlue] at (0.5, 0.5) {$X = [0,1]^2$};
    
    % Point on boundary and its ball
    \coordinate (P) at (1, 0.3);
    \fill[eDarkRed] (P) circle (1.5pt) node[right] {$x$};
    \draw[eDarkRed, dashed] (P) circle (0.3cm);
    
    % Highlighting the part of the ball that is IN X
    \begin{scope}
        \clip (0,0) rectangle (1,1);
        \fill[eDarkRed!40] (P) circle (0.3cm);
    \end{scope}
    
    \node[eDarkRed, align=left, anchor=west] at (1.5, 0.5) {The $\eps$-ball in $X$ is \\ strictly the shaded region! \\ $B_\eps(x) \subseteq X$ is perfectly valid.};
\end{tikzpicture}
\end{center}

\begin{exercise*}
\begin{enumerate}
    \item Write down the definition of convergence of a sequence $(f_n)_{n=0}^\infty$ to $f$ in the metric space $X$ of bounded functions $[0,1] \to \R$ with the supremum metric.
    \item You know the resulting expression from Analysis 1. What is it called?
    \item Use the fact that $C^0([0,1]) \subseteq X$ is closed and the previous proposition of \qt{sequentially closed} to conclude a big result from Analysis 1.
\end{enumerate}
\end{exercise*}

\begin{solution}
\textcolor{eGray}{(Optional)}
\begin{enumerate}
    \item $\forall \eps > 0 \, \exists N \in \N$: $\sup_{x \in [0,1]} \abs{f_n(x) - f(x)} < \eps \quad \forall n \ge N$.
    \item This is \textcolor{eDarkGreen}{uniform convergence}!
    \item $C^0([0,1])$ is closed in the space of bounded functions with the supremum metric. 
    
    \textcolor{eSaddleBrown}{Proof:} Let $f \in X \setminus C^0([0,1])$, i.e. $f:[0,1] \to \R$ is bounded and has at least one discontinuity at a point $x_0 \in [0,1]$. That is, we find $\eps > 0$ s.t. $\forall \delta > 0 \, \exists y \in (x_0 - \delta, x_0 + \delta)$ with $\abs{f(x_0) - f(y)} \ge \eps$.
    
    Then the ball in $X$:
    \[ B_{\eps/3}(f) = \set{g \in X : \sup_{x \in [0,1]} \abs{f(x) - g(x)} < \frac{\eps}{3}} \]
    
    \begin{center}
    \begin{tikzpicture}[scale=1.5, very thick]
        % Axes
        \draw[->, thick, eBlack] (-0.5, 0) -- (4, 0) node[right] {$x$};
        \draw[->, thick, eBlack] (0, -0.5) -- (0, 3) node[above] {$y$};
        \node[below left, eBlack] at (0,0) {$0$};
        \draw[thick, eBlack] (3.5, 0.1) -- (3.5, -0.1) node[below] {$1$};
        
        % Function curves (f)
        \draw[eMediumBlue] (0, 0.5) to[out=0, in=180] (1.5, 0.2);
        \draw[eMediumBlue] (1.5, 1.5) to[out=0, in=180] (3.5, 2.5);
        
        % Discontinuity points
        \fill[eMediumBlue] (1.5, 1.5) circle (1.5pt);
        \fill[eMediumBlue, fill=white, draw=eMediumBlue, thick] (1.5, 0.2) circle (1.5pt);
        
        % Jump epsilon arrow
        \draw[<->, eSaddleBrown] (1.5, 0.3) -- (1.5, 1.4) node[midway, right] {$\eps$};
        
        % Tubes (epsilon/3)
        \draw[eDarkRed, dashed] (0, 0.8) to[out=0, in=180] (1.5, 0.5);
        \draw[eDarkRed, dashed] (0, 0.2) to[out=0, in=180] (1.5, -0.1);
        \draw[eDarkRed, dashed] (1.5, 1.8) to[out=0, in=180] (3.5, 2.8);
        \draw[eDarkRed, dashed] (1.5, 1.2) to[out=0, in=180] (3.5, 2.2);
        
        % Epsilon/3 label arrow
        \draw[<->, eDarkRed] (2.5, 2.05) -- (2.5, 2.35) node[midway, right] {$\frac{\eps}{3}$};
        \node[eMediumBlue] at (3.7, 2.5) {$f$};
    \end{tikzpicture}
    \end{center}
    
    is contained in $X \setminus C^0([0,1])$. Given this fact, the \qt{sequentially closed} proposition ensures that the uniform limit $f$ of $(f_n)_{n=0}^\infty$ in $C^0([0,1])$ is also in $C^0([0,1])$ i.e. continuous, given that $f:[0,1] \to \R$ is bounded.
    
    This is the case, since $f_N$ is bounded, and with $\eps = 1$, we find $N \in \N$ s.t.
    \[ \sup_{x \in [0,1]} \abs{f_N(x) - f(x)} < 1 \implies f \text{ is bounded.} \]
    
    $B_\eps(f)$ in the space of bounded functions on $[0,1]$ is the \qt{$\eps$-tube} around $f$.
\end{enumerate}
\end{solution}



% Week 2 - Friday: Continuity, Compactness
%\section*{Friday}

\begin{warmup}
What is the shape of $B_1(0) \subseteq \R^2$ with:
\begin{itemize}
    \item The Supremum metric?
    \item The Manhattan metric?
\end{itemize}
\end{warmup}

\begin{answer*}\hfill\\ % Forces a line break before the tikzpicture
\begin{center}
\begin{tikzpicture}[scale=1.5, very thick]
    % Supremum Graphic
    \node[eDarkCyan, font=\bfseries] at (0, 1.5) {Supremum:};
    \draw[->, thick, eBlack] (-1.5,0) -- (1.5,0) node[right] {$x_1$};
    \draw[->, thick, eBlack] (0,-1.5) -- (0,1.5) node[above] {$x_2$};
    
    \draw[eDarkGreen, dashed] (-1,-1) rectangle (1,1);
    \fill[eSaddleBrown] (0,1) circle (1.5pt) node[above right] {$(0,1)$};
    \fill[eSaddleBrown] (1,0) circle (1.5pt) node[above right] {$(1,0)$};

    % Manhattan Graphic
    \begin{scope}[xshift=4cm]
        \node[eDarkCyan, font=\bfseries] at (0, 1.5) {Manhattan:};
        \draw[->, thick, eBlack] (-1.5,0) -- (1.5,0) node[right] {$x_1$};
        \draw[->, thick, eBlack] (0,-1.5) -- (0,1.5) node[above] {$x_2$};
        
        \draw[eDarkRed, dashed] (1,0) -- (0,1) -- (-1,0) -- (0,-1) -- cycle;
        \fill[eSaddleBrown] (0,1) circle (1.5pt) node[above right] {$(0,1)$};
        \fill[eSaddleBrown] (1,0) circle (1.5pt) node[above right] {$(1,0)$};
    \end{scope}
\end{tikzpicture}
\end{center}
\end{answer*}

\vspace{0.5cm}

\begin{definition*}[Continuous functions]
For $(X, d_X)$ and $(Y, d_Y)$ metric spaces, a function $f: X \to Y$ is called \emph{continuous}, if one of the following equivalent conditions hold:

\begin{enumerate}
    \item $(\eps-\delta)$: $\forall x_0 \in X, \forall \eps > 0 \, \exists \delta > 0$ s.t. $\forall x \in X$:
    \[ d_X(x_0, x) < \delta \implies d_Y(f(x_0), f(x)) < \eps \]
    
    \item (sequentially continuous):
    \[ \lim_{n \to \infty} f(x_n) = f(x) \]
    For any sequence $(x_n)_{n=0}^\infty$ in $X$ with limit $x \in X$, (and the limit exists).
    
    \item (topological): For all $U \subseteq Y$ open, $f^{-1}(U) \subseteq X$ is \emph{open}.
\end{enumerate}
\end{definition*}

\begin{example*}
Let $f: \R \to \R$, $x \mapsto x^2$. Let $a < b \in \R$.
Then for $a < 0 < b$, $f^{-1}((a,b)) = (-\sqrt{b}, \sqrt{b})$ which is open. \\
For $a < b < 0$, $f^{-1}((a,b)) = \emptyset$. \\
For $0 < a < b$, $f^{-1}((a,b)) = (-\sqrt{b}, -\sqrt{a}) \cup (\sqrt{a}, \sqrt{b})$.

\textcolor{eDarkCyan}{Why is it enough to check only these cases?} \\
\textcolor{eGray}{Hint: $f^{-1}(U_1 \cup U_2) = f^{-1}(U_1) \cup f^{-1}(U_2)$}
\end{example*}

\vspace{0.5cm}

\begin{question}
\begin{enumerate}[start=2] % Matching the numbering from the handwritten notes!
    \item Let $(X, d_X)$ and $(Y, d_Y)$ be metric spaces, where $d_Y$ is the discrete metric. Is $f: (X, d_X) \to (Y, d_Y)$ continuous?
    \item What if $d_X$ is the discrete metric, and $d_Y$ is arbitrary?
\end{enumerate}
\end{question}

\begin{answer}
\begin{enumerate}[start=2]
    \item \textcolor{eDarkRed}{No.} For example, take the spaces $X = Y = \R$, $d_X$ the standard metric, and $f = \text{Id}: X \to Y, z \mapsto z$. Then $\set{0} \subseteq Y$ is open, but $\set{0} = \text{Id}^{-1}(\set{0})$ is not open in $(X, d_X)$.
    \item \textcolor{eDarkGreen}{Yes.} Any subset of $X$ is open in the discrete metric, therefore $f^{-1}(U)$ is open for any open $U \subseteq Y$.
\end{enumerate}
\end{answer}

\hrulefill

\section{Completeness and Compactness}

\begin{definition*}[Cauchy Sequence \& Complete Metric Space]
Let $(X, d)$ be a metric space, and $(x_n)_{n=0}^\infty$ a sequence in $X$. We say that $(x_n)$ is \textcolor{eDarkGreen}{\textbf{Cauchy}}, if:
\[ \forall \eps > 0 \, \exists N \in \N \text{ s.t. } d(x_n, x_m) < \eps \quad \forall n, m \ge N \]

$(X, d)$ is a \textcolor{eDarkGreen}{\textbf{complete metric space}} if every Cauchy sequence in $X$ converges with respect to $d$.
\end{definition*}

\begin{example*}
\begin{itemize}
    \item The Euclidean spaces $(\R^n, \langle \cdot, \cdot \rangle_{\text{eucl}})$ are complete.
    \item Any non-closed subset $U \subseteq \R^n$ (standard metric) and the rationals $\mathbb{Q}$ are not complete.
    \item $(C^0([0,1]), \sup)$ is complete. \textcolor{eGray}{(Zorich, p. 22)}
\end{itemize}
\end{example*}

\begin{definition*}[Compactness]
A metric space $(X, d)$ is compact if every open cover has a finite subcover.
\end{definition*}

\begin{theorem*}
A metric space $(X, d)$ is compact if and only if every sequence in $X$ has a convergent subsequence in $X$.
\end{theorem*}

\begin{center}
\begin{tikzpicture}[scale=1.2, very thick]
    % The metric space X (reshaped so it actually fits inside the cover!)
    \draw[eMediumBlue, fill=eMediumBlue!5] (0.5,0.5) .. controls (-0.5,1.5) and (1,2.8) .. (2.5,2.5) .. controls (4.5,2.2) and (4,0) .. (2,0) .. controls (1,-0.5) and (0.8,0) .. cycle;
    \node[eMediumBlue] at (4.2, 2.5) {\Large $X$};
    
    % U1 covers the bottom left
    \draw[eDarkRed, dashed, fill=eDarkRed!10, fill opacity=0.3] (1.2, 0.4) circle (1.4cm);
    \node[eDarkRed] at (1.2, -0.6) {$U_1$};
    
    % U2 covers the right side
    \draw[eSaddleBrown, dashed, fill=eSaddleBrown!10, fill opacity=0.3] (2.8, 1.2) circle (1.6cm);
    \node[eSaddleBrown] at (3.8, 0.8) {$U_2$};
    
    % U3 covers the top left
    \draw[eMediumOrchid, dashed, fill=eMediumOrchid!10, fill opacity=0.3] (1.0, 1.8) circle (1.3cm);
    \node[eMediumOrchid] at (0.0, 2.6) {$U_3$};
    
    \node[eDarkGray] at (2, -1.5) {$U_1 \cup U_2 \cup U_3 = X$};
\end{tikzpicture}
\end{center}

\textcolor{eGray}{A subset of a metric space $(X, d)$ is compact if the metric space $(K, d)$ is compact. This is equivalent to the definition in the script:}

\vspace{0.5cm}

% --- MAGIC COUNTER TRICK FOR DEF 9.63 ---
\newcounter{tempSection}
\setcounter{tempSection}{\value{section}}
\newcounter{tempTheorem}
\setcounter{tempTheorem}{\value{theorem}}

\setcounter{section}{9}
\setcounter{theorem}{62} % 62 + 1 = 63 when the environment begins

\begin{definition}[Compactness]
Let $(X, d)$ be a metric space. A subset $K \subseteq X$ is called compact if one of the following equivalent conditions hold:
\begin{enumerate}
    \item $K$ is \textcolor{eMediumBlue}{\qt{sequentially compact}}: every sequence $(x_n)_{n \in \N}$ in $K$ has a subsequence that is convergent in $K$.
    \item $K$ is \textcolor{eMediumBlue}{\qt{topologically compact}}: every family of open sets $\set{U_i}_{i \in I}$ that cover $K$, has a finite subcover.
\end{enumerate}
\end{definition}

% --- RESTORE COUNTERS ---
\setcounter{section}{\value{tempSection}}
\setcounter{theorem}{\value{tempTheorem}}

\begin{center}
\begin{tikzpicture}[scale=1.2, thick]
    % Open cover of subset
    \draw[eDarkGray, dashed] (-3, -2) rectangle (5, 3);
    \node[eDarkGray] at (4.5, 2.5) {$X$};

    % The compact set K
    \draw[eDarkRed, fill=eDarkRed!20, rounded corners=15pt] (0,0) -- (1,2) -- (3,1.5) -- (2,-1) -- cycle;
    \node[eDarkRed] at (1.5, 0.5) {\Large $K$};

    % Open sets inside X covering K
    \draw[eMediumBlue, dashed, fill=eMediumBlue!10, fill opacity=0.5, very thick] (0.2, 0.8) circle (1.2cm);
    \draw[eMediumBlue, dashed, fill=eMediumBlue!10, fill opacity=0.5, very thick] (2.2, 0.5) circle (1.3cm);
    \draw[eMediumBlue, dashed, fill=eMediumBlue!10, fill opacity=0.5, very thick] (1.2, -0.2) circle (1cm);
    
    \node[eMediumBlue] at (-0.5, 1.5) {$U_{\alpha_1}$};
    \node[eMediumBlue] at (3.5, 1.2) {$U_{\alpha_2}$};
    \node[eMediumBlue] at (1.2, -1.4) {$U_{\alpha_3}$};
\end{tikzpicture}
\end{center}
\textcolor{eGray}{Note: $U_i \subseteq X$ are open sets in the metric space $(X,d)$.}

\vspace{0.5cm}

\begin{exercises}
\begin{enumerate}
    \item Show that every finite subset $\set{x_1, \dots, x_n} \subseteq X$ is compact.
    \item If $K$ is compact, is it complete?
\end{enumerate}
\end{exercises}

\begin{solution}\hfill\\ % Linebreak for cleaner formatting
\textbf{1.} Let $\set{U_\alpha}_{\alpha \in I}$ be an open cover of $\set{x_1, \dots, x_n}$. For each $i=1, \dots, n$, pick $U_{\alpha_i}$ such that $x_i \in U_{\alpha_i}$. 
Then the collection $\set{U_{\alpha_1}, \dots, U_{\alpha_n}}$ is a finite subcover. Thus, the set is compact.

\vspace{0.3cm}
\textbf{2.} \textcolor{eDarkGreen}{Yes.} Let $(x_n)_{n=0}^\infty$ be a Cauchy sequence in $X$. Then, because $K$ is compact, we find a convergent subsequence:
\[ \lim_{i \to \infty} x_{n_i} = x \in X \]
Fix $\eps > 0$. We find $N_1 \in \N$ such that:
\[ d(x_n, x_m) < \eps \quad \forall n, m > N_1 \]
and $N_2 \in \N$ such that:
\[ d(x, x_{n_i}) < \eps \quad \forall i > N_2 \]
For $N := \max\set{N_1, N_2}$ and $k > N$:
\[ d(x, x_k) \le d(x, x_{n_k}) + d(x_{n_k}, x_k) < 2\eps \]
where $n_k \ge k > N$. So, any Cauchy sequence converges.
\end{solution}
% Week 2
\setcounter{section}{0} 

\section{Structured Spaces}

This semester, you will be introduced to certain structures that can be given to a simple set (which we will often call a \qt{space} $X$ in this context). You can think of this as a hierarchy, where each level grants us more powerful mathematical tools:

\begin{enumerate}
    \item \textbf{Topological spaces:} \textcolor{eGray}{(Focus in the 4th semester)}
    \item \textbf{Metric spaces $(X, d)$:} Defines a mere \emph{distance} between points. Important: $X$ does not necessarily have to be a vector space here! \textcolor{eGray}{(Focus in Analysis)}
    \item \textbf{Normed vector spaces $(X, \norm{\cdot})$:} Defines the \emph{length} of a vector. This presupposes a linear structure (i.e., $X$ must be a vector space). \textcolor{eGray}{(Focus in LinAlg / Analysis)}
    \item \textbf{Inner product spaces $(X, \langle \cdot, \cdot \rangle)$:} Defines both lengths and \emph{angles}! This is the richest structure. \textcolor{eGray}{(Focus in LinAlg)}
\end{enumerate}

\gemininote{Why is this hierarchy important? Because the more powerful structures automatically induce the weaker ones! An inner product always induces a norm via the formula $\norm{x} := \sqrt{\langle x, x \rangle}$. And a norm always induces a metric via the formula $d(x,y) := \norm{x - y}$. The reverse, however, is not generally true!}

\hrulefill

\section{Metric Spaces}

A \textbf{metric space} $(X, d)$ is any set $X$ equipped with a \textbf{metric}. You can think of it as a universal ruler.

Formally, a metric is a function:
\[ d: X \times X \to \mathbb{R}_{\ge 0} \]
such that for all points $x, y, z \in X$, the following three intuitive properties hold:

\begin{enumerate}
    \item \textbf{Definiteness:} $d(x,y) \ge 0$, and $d(x,y) = 0 \iff x=y$. 
    \\ \textcolor{eGray}{In words: \qt{The distance is always positive, and a point has zero distance only to itself.}} 
    
    \item \textbf{Symmetry:} $d(x,y) = d(y,x)$.
    \\ \textcolor{eGray}{In words: \qt{The path from you to me is exactly as long as the path from me to you.}} 
    
    \item \textbf{Triangle Inequality:} $d(x,z) \le d(x,y) + d(y,z)$.
    \\ \textcolor{eGray}{In words: \qt{Taking a detour through a third point can never shorten your trip.}} 
\end{enumerate}

\begin{center}
\begin{tikzpicture}[scale=1.5, very thick]
    \coordinate (X) at (0,0);
    \coordinate (Y) at (2,1.5);
    \coordinate (Z) at (4,0);

    \draw[eMediumBlue, ->] (X) -- node[above left] {$d(x,y)$} (Y);
    \draw[eMediumBlue, ->] (Y) -- node[above right] {$d(y,z)$} (Z);
    \draw[eDarkRed, dashed, ->] (X) -- node[below] {$d(x,z) \le d(x,y) + d(y,z)$} (Z);

    \fill[eSaddleBrown] (X) circle (2pt) node[left] {$x$};
    \fill[eSaddleBrown] (Y) circle (2pt) node[above] {$y$};
    \fill[eSaddleBrown] (Z) circle (2pt) node[right] {$z$};
\end{tikzpicture}
\end{center}

\subsection*{Examples of Metric Spaces}

\begin{example*}[The Discrete Metric]
Let $X$ be any set. The discrete metric is the \qt{digital} approach to measuring distance:
\[
d(x,y) = 
\begin{cases} 
1 & \text{if } x \ne y \\ 
0 & \text{if } x = y 
\end{cases}
\]
\gemininote{This metric is an extremely important tool for constructing counterexamples! It perfectly demonstrates that a metric does not have to match our intuitive geometric understanding of space.}
\end{example*}

\begin{example*}[Metrics in $\R^n$]
We have many choices for measuring distance in $\R^n$. Some common ones are:
\begin{itemize}
    \item \textbf{Manhattan metric ($d_1$):} $d_1(x,y) = \sum_{k=1}^n \abs{x_k - y_k}$ \textcolor{eGray}{(Like a taxi navigating a city grid)}
    \item \textbf{Standard/Euclidean metric ($d_2$):} $d_2(x,y) = \sqrt{\sum_{k=1}^n (x_k - y_k)^2}$ \textcolor{eGray}{(The classical straight-line distance)}
    \item \textcolor{eDarkGreen}{\qt{Supremum metric}} \textbf{($d_3$ or $d_\infty$):} $d_3(x,y) = \sup_{k=1,\dots,n} \abs{x_k - y_k}$ \textcolor{eGray}{(Only the single largest coordinate difference matters)}
\end{itemize}

As a generalization: $d_m(x,y) = \sqrt[m]{\sum_{k=1}^n \abs{x_k - y_k}^m}$ is a valid metric on $\R^n$ for any $m \ge 1$.

\begin{center}
\textbf{Visualizing the \qt{Unit Balls} in $\R^2$:} \\
\vspace{0.3cm}
\begin{tikzpicture}[scale=2]
    \draw[->, thick, eBlack] (-1.5,0) -- (1.5,0) node[right] {$x_1$};
    \draw[->, thick, eBlack] (0,-1.5) -- (0,1.5) node[above] {$x_2$};
    
    \draw[eDarkGreen, very thick, dashed] (-1,-1) rectangle (1,1);
    \node[eDarkGreen] at (1.2, 1.2) {$d_\infty$ (Supremum)};
    
    \draw[eMediumBlue, very thick] (0,0) circle (1cm);
    \node[eMediumBlue] at (1.2, 0.5) {$d_2$ (Standard)};
    
    \draw[eDarkRed, very thick, dotted] (1,0) -- (0,1) -- (-1,0) -- (0,-1) -- cycle;
    \node[eDarkRed] at (0.7, -0.7) {$d_1$ (Manhattan)};
\end{tikzpicture}
\end{center}
\end{example*}

\begin{example*}[Continuous Functions]
Let $X = C^0([0,1], \R)$ denote the set of all continuous functions mapping the interval $[0,1]$ to $\R$. We can equip this function space with the \textcolor{eDarkGreen}{\qt{Supremum Metric}}:
\[ d(f,g) := \sup_{x \in [0,1]} \abs{f(x) - g(x)} \]
\textcolor{eGray}{Intuition: This metric finds the single point $x$ where the two graphs are furthest apart, and defines that maximum vertical gap as the distance between the two functions.}
\end{example*}

\begin{example*}[The n-Sphere]
Let $X = S^2 := \set{ (x_1, x_2, x_3) \in \R^3 : x_1^2 + x_2^2 + x_3^2 = 1 }$ (the 2-dimensional sphere). 
The closest path between two points on the sphere is an arc of a great circle (a geodesic). This defines a unique metric on the surface, similar to how straight lines define a metric on $\R^n$.

\begin{center}
\begin{tikzpicture}[scale=2]
    \draw[thick, eBlack] (0,0) circle (1cm);
    \draw[dashed, eBlack] (1,0) arc (0:180:1cm and 0.3cm);
    \draw[eBlack] (1,0) arc (0:-180:1cm and 0.3cm);
    
    \coordinate (X) at (-0.7, 0.4);
    \coordinate (Y) at (0.6, 0.6);
    
    \fill[eSaddleBrown] (X) circle (1.5pt) node[left] {$x$};
    \fill[eSaddleBrown] (Y) circle (1.5pt) node[right] {$y$};
    
    \draw[eDarkRed, dashed, thick] (X) -- node[below=0.2cm, sloped] {\small Euclidean distance} (Y);
    \draw[eMediumBlue, very thick] (X) to[bend left=20] node[above, sloped] {\small Distance on sphere} (Y);
\end{tikzpicture}
\end{center}
\geminiwarning{Terminology check: A sphere embedded in 3D space is actually a 2-dimensional surface, which is why topologists denote it as $S^2$. The script loosely refers to it as the \qt{n-sphere} here based on the ambient space.}
\end{example*}

\vspace{0.5cm}

\begin{exercise*}
Show that the supremum metric from the continuous functions example is indeed a metric.
\end{exercise*}

\begin{solution}
Symmetry and definiteness follow directly from the properties of the absolute value $\abs{\cdot}$. For the triangle inequality, let's consider three functions $f, g, h \in C^0([0,1])$ and define two helpful auxiliary functions:
\[ p := f - h, \quad q := h - g \]
For any arbitrary but fixed $x \in [0,1]$, the value of the function at that specific point can never exceed the supremum of the entire function:
\[ \abs{p(x)} \le \sup \abs{p} \quad \text{and} \quad \abs{q(x)} \le \sup \abs{q} \]
If we now look at the distance between $f$ and $g$ at the point $x$, we can use the "add zero" trick and apply our inequalities:
\begin{equation}
    \begin{aligned}
        \abs{f(x) - g(x)} &= \abs{p(x) + q(x)} \\
        &\le \abs{p(x)} + \abs{q(x)} \quad \text{\textcolor{eGray}{(Standard triangle inequality in $\R$)}} \\
        &\le \sup \abs{p} + \sup \abs{q} \\
        &= \sup \abs{f - h} + \sup \abs{h - g}
    \end{aligned}
\end{equation}
Because this upper bound holds for \textcolor{eSaddleBrown}{every single} $x \in [0,1]$, it must inherently also hold for the supremum taken over all $x$! Thus:
\[ d(f,g) \le d(f,h) + d(h,g) \]
This proves the triangle inequality.
\end{solution}

\hrulefill

\section{Open and Closed Sets}

Let $(X, d)$ be a metric space, let $x \in X$, and let $r \in \R_{>0}$. First, we define the \textcolor{eSaddleBrown}{\qt{open ball}} of radius $r$:
\[ B_r(x) := \set{ y \in X : d(y,x) < r } \]
\textcolor{eGray}{Note: Pay attention to the strict inequality ($<$). The boundary shell is not part of the ball!}

\begin{definition*}
\begin{itemize}
    \item A subset $U \subseteq X$ is \textcolor{eDarkGreen}{\textbf{open}} iff for \emph{every} $y \in U$, we can find an $\eps > 0$ such that the entire ball $B_\eps(y)$ lies completely inside $U$. \textcolor{eGray}{(You can never accidentally step on the boundary.)}
    \item A subset $A \subseteq X$ is \textcolor{eDarkRed}{\textbf{closed}} iff its complement $X \setminus A$ is open.
\end{itemize}
\end{definition*}

\begin{center}
\begin{tikzpicture}[scale=1.2, very thick]
    \draw[eMediumBlue, fill=eMediumBlue!10, dashed] (0,0) .. controls (1,2) and (3,1) .. (4,0) .. controls (3,-2) and (1,-1) .. (0,0);
    \node[eMediumBlue] at (3.5, 1) {\Large $U$};
    
    \coordinate (Y) at (1.5, 0);
    \draw[eDarkGreen, fill=eDarkGreen!20, dashed] (Y) circle (0.6cm);
    \fill[eSaddleBrown] (Y) circle (2pt) node[below] {$y$};
    \draw[eDarkGreen, ->] (Y) -- ++(45:0.6cm) node[midway, above left] {$\eps$};
    \node[eDarkGreen] at (2.5, -0.5) {$B_\eps(y) \subseteq U$};
\end{tikzpicture}
\end{center}

\vspace{0.5cm}

\begin{exercise*}
Let $(X,d)$ be a metric space. Show that the entire space $X$ and the empty set $\emptyset$ are both open and closed.
\end{exercise*}

\begin{solution}
First, let's look at openness:
\begin{itemize}
    \item $X$ is open because for any $x \in X$ and any radius $r > 0$, the ball $B_r(x)$ is by definition a subset of $X$. Thus, $B_r(x) \subseteq X$ trivially holds.
    \item $\emptyset$ contains no points. Therefore, the condition \qt{for every point $y \in \emptyset$ there exists...} is a vacuously true statement. Thus, the empty set is open.
\end{itemize}
Now for closedness:
Since $X$ and $\emptyset$ are each other's complements in $X$ (i.e., $X \setminus X = \emptyset$ and $X \setminus \emptyset = X$), and we just proved both are open, they must both inherently be closed as well!
\end{solution}

\hrulefill

\section{Sequences and Convergence}

A sequence $(x_n)_{n=0}^\infty$ in a metric space $(X, d)$ is formally a function from $\N \to X$. We say that $(x_n)$ \textbf{converges} to a limit $x \in X$ if:
\[ \forall \eps > 0, \exists N \in \N \text{ such that } d(x_n, x) < \eps \quad \forall n \ge N \]

\subsection*{The Big Connection Proposition}
Topological properties (openness/closedness) and the behavior of sequences go hand-in-hand:
\begin{enumerate}
    \item $U \subseteq X$ is \textbf{open} $\iff$ for every sequence in $X$ that has a limit in $U$, the sequence eventually lies entirely inside $U$.
    \item $A \subseteq X$ is \textbf{closed} $\iff$ \emph{every} convergent sequence that lies entirely in $A$ must necessarily have its limit in $A$ as well. \textcolor{eGray}{(The sequence cannot \qt{break out} of the set.)}
\end{enumerate}

\subsection*{\warning{Watch out! A Common Misconception!}}
Let's look at the standard Euclidean metric space $\R^2$. We know the square $[0,1]^2$ is closed and $(0,1)^2$ is open. (The rectangle $[0,1] \times (0,1)$ is neither open nor closed in $\R^2$).

\textcolor{eDarkRed}{But...} what if our \textit{entire universe} (our metric space $X$) is defined strictly as $X := [0,1]^2$? 
In the subspace $X = [0,1]^2$, the set $[0,1]^2$ is suddenly \emph{open}! 

\textcolor{eDarkCyan}{Why?} Because the full space $X$ is always open within itself. If you take a point exactly on the boundary, its open $\eps$-ball consists only of points \textcolor{eSaddleBrown}{that actually exist in $X$}. The ball does not bleed out into $\R^2$ because $\R^2$ simply does not exist for this specific metric space!

\begin{center}
\begin{tikzpicture}[scale=2, very thick]
    \draw[->, eBlack] (-0.5, 0) -- (1.5, 0) node[right] {$x_1$};
    \draw[->, eBlack] (0, -0.5) -- (0, 1.5) node[above] {$x_2$};
    
    \draw[eMediumBlue, fill=eMediumBlue!10] (0,0) rectangle (1,1);
    \node[eMediumBlue] at (0.5, 0.5) {$X = [0,1]^2$};
    
    \coordinate (P) at (1, 0.3);
    \fill[eDarkRed] (P) circle (1.5pt) node[right] {$x$};
    \draw[eDarkRed, dashed] (P) circle (0.3cm);
    
    \begin{scope}
        \clip (0,0) rectangle (1,1);
        \fill[eDarkRed!40] (P) circle (0.3cm);
    \end{scope}
    
    \node[eDarkRed, align=left, anchor=west] at (1.5, 0.5) {The $\eps$-ball in $X$ is \\ strictly the shaded region! \\ $B_\eps(x) \subseteq X$ is perfectly valid.};
\end{tikzpicture}
\end{center}

\begin{exercise*}
\begin{enumerate}
    \item Write down the definition of convergence for a sequence of functions $(f_n)_{n=0}^\infty$ to $f$ in the metric space $X$ of bounded functions $[0,1] \to \R$ equipped with the supremum metric.
    \item You know this resulting expression from Analysis 1. What kind of convergence is this called?
    \item Use the fact that the space of continuous functions $C^0([0,1]) \subseteq X$ is closed, and combine it with the \qt{sequentially closed} proposition (part b) to conclude a massive result from Analysis 1.
\end{enumerate}
\end{exercise*}

\begin{solution}
\begin{enumerate}
    \item $\forall \eps > 0 \, \exists N \in \N$ such that $\sup_{x \in [0,1]} \abs{f_n(x) - f(x)} < \eps \quad \forall n \ge N$.
    \item This is exactly the definition of \textcolor{eDarkGreen}{\textbf{Uniform Convergence}}!
    \item The fact that $C^0([0,1])$ is closed in the space of bounded functions means (according to our proposition above): If a sequence of continuous functions converges uniformly to a limit function $f$, then this limit function $f$ cannot leave the space of continuous functions. Ergo: \textbf{The uniform limit of continuous functions is always continuous!}
    
    \textcolor{eSaddleBrown}{Proof of closedness (The $\eps/3$-trick):} 
    We show that the complement is open. Let $f \in X \setminus C^0([0,1])$. This means $f$ is bounded but has at least one discontinuity at a point $x_0$. That implies there is a real jump $\eps > 0$, such that no matter how close ($\delta$) we get to $x_0$, we can always find a $y$ where $\abs{f(x_0) - f(y)} \ge \eps$.
    
    Now we consider an open ball around this discontinuous function $f$, but we give it a tiny radius of $\eps/3$:
    \[ B_{\eps/3}(f) = \set{g \in X : \sup_{x \in [0,1]} \abs{f(x) - g(x)} < \frac{\eps}{3}} \]
    
    \begin{center}
    \begin{tikzpicture}[scale=1.5, very thick]
        \draw[->, thick, eBlack] (-0.5, 0) -- (4, 0) node[right] {$x$};
        \draw[->, thick, eBlack] (0, -0.5) -- (0, 3) node[above] {$y$};
        \node[below left, eBlack] at (0,0) {$0$};
        \draw[thick, eBlack] (3.5, 0.1) -- (3.5, -0.1) node[below] {$1$};
        
        \draw[eMediumBlue] (0, 0.5) to[out=0, in=180] (1.5, 0.2);
        \draw[eMediumBlue] (1.5, 1.5) to[out=0, in=180] (3.5, 2.5);
        
        \fill[eMediumBlue] (1.5, 1.5) circle (1.5pt);
        \fill[eMediumBlue, fill=white, draw=eMediumBlue, thick] (1.5, 0.2) circle (1.5pt);
        
        \draw[<->, eSaddleBrown] (1.5, 0.3) -- (1.5, 1.4) node[midway, right] {Jump $\eps$};
        
        \draw[eDarkRed, dashed] (0, 0.8) to[out=0, in=180] (1.5, 0.5);
        \draw[eDarkRed, dashed] (0, 0.2) to[out=0, in=180] (1.5, -0.1);
        \draw[eDarkRed, dashed] (1.5, 1.8) to[out=0, in=180] (3.5, 2.8);
        \draw[eDarkRed, dashed] (1.5, 1.2) to[out=0, in=180] (3.5, 2.2);
        
        \draw[<->, eDarkRed] (2.5, 2.05) -- (2.5, 2.35) node[midway, right] {$\frac{\eps}{3}$};
        \node[eMediumBlue] at (3.7, 2.5) {$f$};
    \end{tikzpicture}
    \end{center}
    
    This ball (the \qt{$\eps/3$-tube} around $f$) contains exclusively functions that are forced to replicate this jump. No function $g$ trapped inside this tube can possibly be continuous! Therefore, the entire ball $B_{\eps/3}(f)$ is contained within the complement $X \setminus C^0([0,1])$. Since the complement is open, $C^0([0,1])$ is closed!
\end{enumerate}
\end{solution}

\hrulefill

\section*{Friday}

\begin{warmup}
What is the shape of the open unit ball $B_1(0) \subseteq \R^2$ with:
\begin{itemize}
    \item The Supremum metric?
    \item The Manhattan metric?
\end{itemize}
\end{warmup}

\begin{answer*}\hfill\\ 
\begin{center}
\begin{tikzpicture}[scale=1.5, very thick]
    \node[eDarkCyan, font=\bfseries] at (0, 1.5) {Supremum ($d_\infty$):};
    \draw[->, thick, eBlack] (-1.5,0) -- (1.5,0) node[right] {$x_1$};
    \draw[->, thick, eBlack] (0,-1.5) -- (0,1.5) node[above] {$x_2$};
    
    \draw[eDarkGreen, dashed] (-1,-1) rectangle (1,1);
    \fill[eSaddleBrown] (0,1) circle (1.5pt) node[above right] {$(0,1)$};
    \fill[eSaddleBrown] (1,0) circle (1.5pt) node[above right] {$(1,0)$};

    \begin{scope}[xshift=4cm]
        \node[eDarkCyan, font=\bfseries] at (0, 1.5) {Manhattan ($d_1$):};
        \draw[->, thick, eBlack] (-1.5,0) -- (1.5,0) node[right] {$x_1$};
        \draw[->, thick, eBlack] (0,-1.5) -- (0,1.5) node[above] {$x_2$};
        
        \draw[eDarkRed, dashed] (1,0) -- (0,1) -- (-1,0) -- (0,-1) -- cycle;
        \fill[eSaddleBrown] (0,1) circle (1.5pt) node[above right] {$(0,1)$};
        \fill[eSaddleBrown] (1,0) circle (1.5pt) node[above right] {$(1,0)$};
    \end{scope}
\end{tikzpicture}
\end{center}
\end{answer*}

\vspace{0.5cm}

\begin{definition*}[Continuous functions]
For two metric spaces $(X, d_X)$ and $(Y, d_Y)$, a function $f: X \to Y$ is called \emph{continuous} if one of the following, completely equivalent, conditions holds:

\begin{enumerate}
    \item \textbf{$(\eps-\delta)$-criterion:} $\forall x_0 \in X, \forall \eps > 0 \, \exists \delta > 0$ such that $\forall x \in X$:
    \[ d_X(x_0, x) < \delta \implies d_Y(f(x_0), f(x)) < \eps \]
    \textcolor{eGray}{(If you are close enough to $x_0$ in the starting space, you are guaranteed to land close to $f(x_0)$ in the target space.)}
    
    \item \textbf{Sequential continuity:} For any sequence $(x_n)_{n=0}^\infty$ in $X$ with a limit $x \in X$, it holds that:
    \[ \lim_{n \to \infty} f(x_n) = f(x) \]
    \textcolor{eGray}{(Continuous functions allow you to pull the limit inside the argument.)}
    
    \item \textbf{Topological continuity:} For all open sets $U \subseteq Y$, the preimage $f^{-1}(U) \subseteq X$ is \emph{open}.
    \textcolor{eGray}{(Continuous functions preserve the topological structure when mapping backwards.)}
\end{enumerate}
\end{definition*}

\begin{example*}
Let $f: \R \to \R$, $x \mapsto x^2$ and let $(a,b) \subseteq \R$ be an interval. The preimage of this interval depends heavily on $a$ and $b$:
\begin{itemize}
    \item For $a < 0 < b$: $f^{-1}((a,b)) = (-\sqrt{b}, \sqrt{b})$. \textcolor{eGray}{(This is an open interval!)}
    \item For $a < b < 0$: $f^{-1}((a,b)) = \emptyset$. \textcolor{eGray}{(The empty set is, as we proved, open!)}
    \item For $0 < a < b$: $f^{-1}((a,b)) = (-\sqrt{b}, -\sqrt{a}) \cup (\sqrt{a}, \sqrt{b})$. \textcolor{eGray}{(A union of open intervals is open!)}
\end{itemize}
\textcolor{eDarkCyan}{Why is it enough to check only these simple basis intervals?} \\
\textcolor{eGray}{Hint: The preimage perfectly respects unions: $f^{-1}(U_1 \cup U_2) = f^{-1}(U_1) \cup f^{-1}(U_2)$. Any open set on $\R$ can be constructed by patching together such intervals.}
\end{example*}

\vspace{0.5cm}

\begin{question*}
Consider functions mapping between two metric spaces $(X, d_X)$ and $(Y, d_Y)$.
\begin{enumerate}
    \item What happens if $d_Y$ is the discrete metric? Is a standard function like $f: (\R, d_{std}) \to (\R, d_{discrete})$ continuous?
    \item And conversely: What if the starting space $d_X$ is discrete, and the target space $d_Y$ is entirely arbitrary?
\end{enumerate}
\end{question*}

\begin{answer*}
\begin{enumerate}
    \item \textcolor{eDarkRed}{No, almost never.} Take the identity function $f(z) = z$ as an example. In the discrete space $Y$, even a single point like $\set{0}$ is an open set! Let's look at its preimage: $\text{Id}^{-1}(\set{0}) = \set{0}$. But a single point is definitely \textbf{not} open in the starting space with the standard metric $d_{std}$. The preimage of an open set is not open here, so the function is not continuous!
    
    \item \textcolor{eDarkGreen}{Yes, always!} In a discrete space, every single subset is open. This means that no matter what open set $U \subseteq Y$ you pick, its preimage $f^{-1}(U)$ is guaranteed to be open in $X$ (because \emph{everything} is open there). Any function mapping out of a discrete space is automatically continuous.
\end{enumerate}
\end{answer*}

\hrulefill

\section{Completeness and Compactness}

\begin{definition*}[Cauchy Sequences \& Complete Metric Spaces]
Let $(X, d)$ be a metric space, and $(x_n)_{n=0}^\infty$ a sequence in $X$. We call $(x_n)$ a \textcolor{eDarkGreen}{\textbf{Cauchy sequence}} if its terms get arbitrarily close to one another beyond a certain point:
\[ \forall \eps > 0 \, \exists N \in \N \text{ such that } d(x_n, x_m) < \eps \quad \forall n, m \ge N \]

A metric space $(X, d)$ is called a \textcolor{eDarkGreen}{\textbf{complete metric space}} if every Cauchy sequence in $X$ actually converges to a real limit within $X$.
\end{definition*}

\gemininote{Intuition for completeness: A Cauchy sequence desperately "wants" to converge. A space is incomplete if it sabotages this—for instance, if the space has \qt{holes} (like $\mathbb{Q}$) or is missing its boundary, meaning the sequence steers towards a limit that isn't actually there.}

\begin{example*}
\begin{itemize}
    \item The Euclidean spaces $(\R^n, \langle \cdot, \cdot \rangle_{\text{eucl}})$ are complete.
    \item Any non-closed subset $U \subseteq \R^n$ and the rational numbers $\mathbb{Q}$ are \emph{not} complete.
    \item The space of continuous functions $(C^0([0,1]), \sup)$ is complete!
\end{itemize}
\end{example*}

\begin{definition*}[Compactness]
A metric space $(X, d)$ is compact if every wild, infinite open cover can be reduced to a \emph{finite} subcover.
\end{definition*}

\begin{theorem*}
A metric space $(X, d)$ is compact if and only if \emph{every} sequence in $X$ has a subsequence that converges to a point in $X$.
\end{theorem*}

\begin{center}
\begin{tikzpicture}[scale=1.2, very thick]
    \draw[eMediumBlue, fill=eMediumBlue!5] (0.5,0.5) .. controls (-0.5,1.5) and (1,2.8) .. (2.5,2.5) .. controls (4.5,2.2) and (4,0) .. (2,0) .. controls (1,-0.5) and (0.8,0) .. cycle;
    \node[eMediumBlue] at (4.2, 2.5) {\Large $X$};
    
    \draw[eDarkRed, dashed, fill=eDarkRed!10, fill opacity=0.3] (1.2, 0.4) circle (1.4cm);
    \node[eDarkRed] at (1.2, -0.6) {$U_1$};
    
    \draw[eSaddleBrown, dashed, fill=eSaddleBrown!10, fill opacity=0.3] (2.8, 1.2) circle (1.6cm);
    \node[eSaddleBrown] at (3.8, 0.8) {$U_2$};
    
    \draw[eMediumOrchid, dashed, fill=eMediumOrchid!10, fill opacity=0.3] (1.0, 1.8) circle (1.3cm);
    \node[eMediumOrchid] at (0.0, 2.6) {$U_3$};
    
    \node[eDarkGray] at (2, -1.5) {$U_1 \cup U_2 \cup U_3 = X$};
\end{tikzpicture}
\end{center}

\textcolor{eGray}{Note: A subset $K$ of a metric space $(X, d)$ is compact if the subspace $(K, d)$ is compact. This is strictly equivalent to the following prominent definition from the official script:}

\vspace{0.5cm}

% --- MAGIC COUNTER TRICK FOR DEF 9.63 ---
\newcounter{tempSection}
\setcounter{tempSection}{\value{section}}
\newcounter{tempTheorem}
\setcounter{tempTheorem}{\value{theorem}}

\setcounter{section}{9}
\setcounter{theorem}{62} 

\begin{definition}[Compactness]
Let $(X, d)$ be a metric space. A subset $K \subseteq X$ is called compact if one (and therefore both) of the following equivalent conditions hold:
\begin{enumerate}
    \item $K$ is \textcolor{eMediumBlue}{\qt{sequentially compact}}: every sequence $(x_n)_{n \in \N}$ in $K$ has a subsequence that is convergent in $K$.
    \item $K$ is \textcolor{eMediumBlue}{\qt{topologically compact}}: every family of open sets $\set{U_i}_{i \in I}$ that cover $K$ has a \emph{finite} subcover.
\end{enumerate}
\end{definition}

% --- RESTORE COUNTERS ---
\setcounter{section}{\value{tempSection}}
\setcounter{theorem}{\value{tempTheorem}}

\begin{center}
\begin{tikzpicture}[scale=1.2, thick]
    \draw[eDarkGray, dashed] (-3, -2) rectangle (5, 3);
    \node[eDarkGray] at (4.5, 2.5) {$X$};

    \draw[eDarkRed, fill=eDarkRed!20, rounded corners=15pt] (0,0) -- (1,2) -- (3,1.5) -- (2,-1) -- cycle;
    \node[eDarkRed] at (1.5, 0.5) {\Large $K$};

    \draw[eMediumBlue, dashed, fill=eMediumBlue!10, fill opacity=0.5, very thick] (0.2, 0.8) circle (1.2cm);
    \draw[eMediumBlue, dashed, fill=eMediumBlue!10, fill opacity=0.5, very thick] (2.2, 0.5) circle (1.3cm);
    \draw[eMediumBlue, dashed, fill=eMediumBlue!10, fill opacity=0.5, very thick] (1.2, -0.2) circle (1cm);
    
    \node[eMediumBlue] at (-0.5, 1.5) {$U_{\alpha_1}$};
    \node[eMediumBlue] at (3.5, 1.2) {$U_{\alpha_2}$};
    \node[eMediumBlue] at (1.2, -1.4) {$U_{\alpha_3}$};
\end{tikzpicture}
\end{center}
\textcolor{eGray}{Important: The sets $U_i \subseteq X$ in the cover are always \emph{open sets} with respect to the ambient space $(X,d)$.}

\vspace{0.5cm}

\begin{exercises}
\begin{enumerate}
    \item Show that every finite subset $\set{x_1, \dots, x_n} \subseteq X$ is compact.
    \item If a set $K$ is compact, is it automatically complete?
\end{enumerate}
\end{exercises}

\begin{solution}\hfill\\ 
\textbf{a)} Let's use topological compactness: Let $\set{U_\alpha}_{\alpha \in I}$ be an arbitrary open cover of $\set{x_1, \dots, x_n}$. Since these sets must cover all our points, we can simply pick \emph{one} set $U_{\alpha_i}$ from the massive pool for each point $x_i$ (from $i=1$ to $n$) such that $x_i \in U_{\alpha_i}$. 
The collection $\set{U_{\alpha_1}, \dots, U_{\alpha_n}}$ is then a subcover. Because it consists of only $n$ (a finite number!) sets, we have found a finite subcover. Thus, the finite set is compact.

\vspace{0.3cm}
\textbf{b)} \textcolor{eDarkGreen}{Yes, absolutely.} Sequential compactness is perfect for this: Let $(x_n)_{n=0}^\infty$ be an arbitrary Cauchy sequence in $X$. Because $K$ is (sequentially) compact, this sequence \emph{must} have a convergent subsequence:
\[ \lim_{i \to \infty} x_{n_i} = x \in X \]
Now a fundamental trick of Cauchy sequences comes into play: If a Cauchy sequence agrees to converge even for a fraction of its terms (the subsequence), it forcibly drags the rest of the sequence along with it!

Formally: Fix an $\eps > 0$. Because it is a Cauchy sequence, all terms move within $\eps$-distance of each other beyond some index $N_1$:
\[ d(x_n, x_m) < \eps \quad \forall n, m > N_1 \]
Because the subsequence converges to $x$, its terms move within $\eps$-distance of $x$ beyond some index $N_2$:
\[ d(x, x_{n_i}) < \eps \quad \forall i > N_2 \]
If we go far enough into the sequence (choosing $N := \max\set{N_1, N_2}$), we can estimate the distance of any arbitrary term $x_k$ (with $k > N$) to the limit $x$. We just cleverly insert an element $x_{n_k}$ of the convergent subsequence in between:
\[ d(x, x_k) \le d(x, x_{n_k}) + d(x_{n_k}, x_k) < 2\eps \]
\textcolor{eGray}{(Where of course $n_k \ge k > N$).} Since we can make the distance $2\eps$ arbitrarily small, the entire Cauchy sequence indeed converges to $x$. Thus, every compact space is always complete!
\end{solution}

% Week 3
\setcounter{section}{2} % Set to appropriate section number for Week 3

\section{Compactness (Continued)}

\begin{theorem*}[Heine-Borel]
A subset $K \subseteq \R^n$ equipped with the \emph{standard metric} is compact if and only if it is closed and bounded.
\end{theorem*}

\gemininote{The Heine-Borel theorem perfectly captures our intuitive visual understanding of a "compact" (small, contained, boundary-inclusive) set in standard geometry. However, this beautiful equivalence strictly relies on the standard Euclidean space. We will now prove that Heine-Borel completely breaks down if we use a non-standard metric!}

\begin{exercise*}
Let $(X, d)$ be a metric space.
\begin{enumerate}
    \item Show that $\tilde{d}(x,y) := \frac{d(x,y)}{1 + d(x,y)}$ is also a valid metric on $X$.
    \item Show that if $U \subseteq X$ is open with respect to $d$, then $U$ is also open with respect to $\tilde{d}$. \textcolor{eGray}{(Optional)}
    \item Show that for $X = \R^n$ with the metric $\tilde{d}(x,y) = \frac{\abs{x-y}}{1+\abs{x-y}}$, the entire space $\R^n \subseteq X$ is closed and bounded, but \emph{not} compact.
\end{enumerate}
\end{exercise*}

\begin{solution}
\textbf{a)} Non-negativity, definiteness, and symmetry follow immediately from the properties of $d(x,y)$. 
For the triangle inequality, we want to show:
\[ \frac{d(x,y)}{1+d(x,y)} \le \frac{d(x,z)}{1+d(x,z)} + \frac{d(z,y)}{1+d(z,y)} \]
Let's denote $d_1 = d(x,y)$, $d_2 = d(x,z)$, and $d_3 = d(z,y)$. We cross-multiply to clear the denominators:
\[ d_1(1+d_2)(1+d_3) \le d_2(1+d_1)(1+d_3) + d_3(1+d_1)(1+d_2) \]
Expanding everything yields:
\[ d_1 + d_1 d_2 + d_1 d_3 + d_1 d_2 d_3 \le d_2 + d_1 d_2 + d_2 d_3 + d_1 d_2 d_3 + d_3 + d_1 d_3 + d_2 d_3 + d_1 d_2 d_3 \]
Canceling the common terms ($d_1 d_2$, $d_1 d_3$, and one $d_1 d_2 d_3$) leaves us with:
\[ d_1 \le d_2 + 2d_2 d_3 + d_3 + d_1 d_2 d_3 \]
We know from the standard triangle inequality that $d_1 \le d_2 + d_3$. Subtracting this known truth from both sides yields:
\[ 0 \le 2d_2 d_3 + d_1 d_2 d_3 \]
Since distances are non-negative, this final inequality is trivially true! Thus, $\tilde{d}$ is a metric.

\vspace{0.3cm}
\textbf{b)} \textcolor{eGray}{(Proof sketch)} The function $t \mapsto \frac{t}{1+t}$ is strictly monotonically increasing for $t \ge 0$. Therefore, open balls in $d$ can be directly mapped to open balls in $\tilde{d}$ by simply transforming the radius $r$ into $\tilde{r} = \frac{r}{1+r}$.

\vspace{0.3cm}
\textbf{c)} Note that for any $t \ge 0$, the expression $\frac{t}{1+t}$ is strictly strictly less than $1$. Therefore, for all $x,y \in \R^n$:
\[ \tilde{d}(x,y) = \frac{\abs{x-y}}{1+\abs{x-y}} < 1 \]
This means the entire space $\R^n$ fits inside the open ball $\tilde{B}_1(0)$. Thus, $\R^n$ is \textbf{bounded} with respect to this metric! Also, the full space is trivially \textbf{closed}. 
However, consider the sequence $x_n = (n, 0, \dots, 0)$. This sequence wanders off to infinity and has absolutely no convergent subsequence. Since it lacks sequential compactness, $\R^n$ is \textbf{not compact}, breaking Heine-Borel!
\end{solution}

\hrulefill

\section{Why do we care about compactness?}

Compactness is one of the most powerful properties a space can have. It grants us three massive mathematical guarantees:

\begin{enumerate}
    \item \textbf{Preserved by continuous functions:} 
    Let $X, Y$ be metric spaces, and $f: X \to Y$ continuous. If $K \subseteq X$ is compact, then its image $f(K) \subseteq Y$ is guaranteed to be compact.
    
    \item \textbf{Weierstrass Theorem / Extreme Value Theorem:}
    Let $K \subseteq X$ be a compact metric space, and $f: K \to \R$ a continuous function. Then $f$ attains an absolute maximum and minimum within $K$. \textcolor{eGray}{(In particular, $f$ is bounded!)}
    
    \item \textbf{Uniform continuity:}
    Let $X, Y$ be metric spaces, and $f: X \to Y$ continuous. If $X$ is compact, then $f$ is automatically \emph{uniformly} continuous! \textcolor{eGray}{(You can pick a global $\delta$ that works everywhere for a given $\eps$.)}
\end{enumerate}

\begin{center}
\begin{tikzpicture}[scale=1.5, thick]
    % Axes
    \draw[->, eBlack] (-0.5, 0) -- (4, 0) node[right] {$x$};
    \draw[->, eBlack] (0, -0.5) -- (0, 3) node[above] {$f(x)$};
    
    % Function curve
    \draw[eMediumBlue, very thick] (0.5, 1.5) .. controls (1, 2.5) and (1.5, 3) .. (2.5, 0.5) .. controls (3, 0) and (3.5, 1) .. (3.8, 1.5);
    
    % Interval markers
    \draw[dashed, eGray] (0.5, 0) node[below, eBlack] {$a$} -- (0.5, 1.5);
    \draw[dashed, eGray] (3.8, 0) node[below, eBlack] {$b$} -- (3.8, 1.5);
    
    % Max / Min markers
    \coordinate (Max) at (1.15, 2.55);
    \coordinate (Min) at (2.7, 0.35);
    
    \fill[eDarkRed] (Max) circle (2pt);
    \fill[eDarkCyan] (Min) circle (2pt);
    
    \draw[dashed, eGray] (Max) -- (1.15, 0) node[below, eBlack] {$c$};
    \draw[dashed, eGray] (Max) -- (0, 2.55) node[left, eBlack] {$f(c)$};
    
    \draw[dashed, eGray] (Min) -- (2.7, 0) node[below, eBlack] {$d$};
    \draw[dashed, eGray] (Min) -- (0, 0.35) node[left, eBlack] {$f(d)$};
\end{tikzpicture}\\
\textcolor{eGray}{A continuous function on a compact interval $[a,b]$ always attains its absolute maximum (red) and absolute minimum (cyan).}
\end{center}

\hrulefill

\section{Connectedness}

\begin{definition*}[Disconnected metric space/subset]
A metric space $(X, d)$ is \textbf{disconnected} if there exist subsets $U_1, U_2 \subseteq X$ which are:
\begin{enumerate}
    \item \textbf{open}
    \item \textbf{non-empty} ($U_1 \ne \emptyset \ne U_2$)
    \item \textbf{disjoint} ($U_1 \cap U_2 = \emptyset$)
\end{enumerate}
such that $X = U_1 \cup U_2$. Similarly, a subset $A \subseteq X$ is disconnected if it can be split in this exact same way using the induced metric. \\
\textbf{Connected} is simply defined as the logical negation of disconnected!
\end{definition*}

\gemininote{Intuitively, a space is disconnected if it consists of two or more completely isolated \qt{islands}. Because proving connectedness directly is very hard, we almost always prove it by contradiction: We assume the space \emph{is} disconnected (split into two islands $U_1$ and $U_2$) and show that this leads to a logical paradox.}

\begin{center}
\begin{tikzpicture}[scale=1.2, very thick]
    % Disconnected
    \node[eDarkRed, font=\bfseries] at (1.5, 1.5) {Disconnected};
    \draw[eMediumBlue, fill=eMediumBlue!10] (0,0) circle (0.8cm);
    \draw[eMediumBlue, fill=eMediumBlue!10] (3,0) circle (1cm);
    \node[eMediumBlue] at (0,0) {$U_1$};
    \node[eMediumBlue] at (3,0) {$U_2$};
    \draw[<->, eDarkRed, dashed] (0.9,0) -- (1.9,0) node[midway, above] {Gap};

    % Connected
    \begin{scope}[xshift=6cm]
        \node[eDarkGreen, font=\bfseries] at (1.5, 1.5) {Connected};
        \draw[eMediumBlue, fill=eMediumBlue!10] (0,0) .. controls (1,1) and (2,1) .. (3,0) .. controls (2,-1) and (1,-1) .. (0,0);
        \node[eMediumBlue] at (1.5,0) {$X$};
    \end{scope}
\end{tikzpicture}
\end{center}

\begin{theorem*}
Continuous functions preserve connected sets! Let $X, Y$ be metric spaces, and $f: X \to Y$ continuous. If $A \subseteq X$ is connected, then its image $f(A)$ is strictly connected in $Y$.
\end{theorem*}

\begin{exercise*}
Are the following statements True or False?
\begin{enumerate}
    \item The union of connected subsets with a common point is connected.
    \item The intersection of connected subsets is connected.
    \item The closure of a connected subset is connected.
\end{enumerate}
\end{exercise*}

\begin{solution}
\textbf{1. True.} \\
\textcolor{eSaddleBrown}{Proof:} Let $X$ be a metric space, $V_1, V_2 \subseteq X$ connected, and let $x_0 \in V_1 \cap V_2$ be their shared point. Assume by contradiction that their union is disconnected. Then $V_1 \cup V_2 = U_1 \cup U_2$, where $U_1, U_2 \subseteq X$ are open, non-empty, and disjoint. 
Without loss of generality, assume the shared point $x_0$ lies in $U_1$. Because $x_0$ is in both $V_1$ and $V_2$, it follows that $U_1 \cap V_1 \ne \emptyset$ and $U_1 \cap V_2 \ne \emptyset$. 
Since $V_1$ and $V_2$ are individually connected, they cannot be split across $U_1$ and $U_2$. Since they both share a point with $U_1$, they must both lie \emph{entirely} inside $U_1$. Thus $V_1 \cap U_2 = \emptyset$ and $V_2 \cap U_2 = \emptyset$. This forces $U_2 = \emptyset$, which violates the definition of disconnectedness! Contradiction.

\vspace{0.3cm}
\textbf{2. False.} \\
\textcolor{eSaddleBrown}{Counterexample:} Let $X = \R^2$. Let $U = [-1, 1] \times \set{0}$ (a horizontal line segment) and $V = S^1$ (the unit circle). Both are perfectly connected. However, their intersection is exactly two isolated points: $U \cap V = \set{(-1,0), (1,0)}$, which is clearly disconnected!

\vspace{0.3cm}
\textbf{3. True.} \\
\textcolor{eSaddleBrown}{Proof:} Let $A$ be connected with closure $\bar{A}$. Suppose $\bar{A}$ is disconnected, meaning there exist disjoint open sets $U_1, U_2$ covering $\bar{A}$. Without loss of generality, let $A \cap U_1 \ne \emptyset$. Since $A$ is connected and covered by $U_1 \cup U_2$, it must lie entirely in one of them, meaning $A \cap U_2 = \emptyset$. 
However, $U_2$ is an open set. If an open set contains no points of $A$, it cannot possibly contain any boundary points of $A$ either! Thus $U_2 \cap \bar{A} = \emptyset$, meaning $U_2$ is empty. Contradiction.
\end{solution}

\hrulefill

\section{Path Connectedness}

\begin{definition*}[Paths and Path-connectedness]
Let $(X, d)$ be a metric space. 
\begin{itemize}
    \item A \textbf{path} is a continuous function $\gamma: [0,1] \to X$, where $[0,1]$ carries the standard metric.
    \item A space (or subset) $X$ is \textbf{path-connected} if for any two points $x, y \in X$, there exists a path $\gamma$ such that $\gamma(0) = x$ and $\gamma(1) = y$.
\end{itemize}
\end{definition*}

\begin{proposition*}
\begin{itemize}
    \item \textbf{Path-connected $\implies$ connected.} (This is always true!)
    \item For \emph{open} subsets of Euclidean space $\R^n$: \textbf{connected $\iff$ path-connected.}
\end{itemize}
\end{proposition*}

\subsection*{\warning{Counterexample: The Topologist's Sine Curve}}
The reverse implication (connected $\implies$ path-connected) is \textbf{false} for non-open sets! Consider the set:
\[ S := \set{\left(t, \sin\left(\frac{1}{t}\right)\right) : t > 0} \cup \set{(0,0)} \subseteq \R^2 \]
$S$ is \textbf{connected} (because the sine wave part is connected, and we proved earlier that adding boundary points via the closure maintains connectedness). However, $S$ is \textbf{not path-connected}!

\gemininote{Why can't we draw a path? As $t$ approaches $0$, the frequency of the sine wave goes to infinity. It oscillates up and down infinitely fast between $-1$ and $1$. If you try to draw a path from the origin $(0,0)$ to the rest of the curve, your pen would have to cross the x-axis an infinite number of times in a finite amount of time, shattering continuity!}

\begin{solution}
Assume by contradiction that $S$ is path-connected. Then there exists a continuous path $\gamma: [0,1] \to S$ starting at the origin $\gamma(0) = (0,0)$ and ending at some point on the curve, say $\gamma(1) = (1, \sin(1))$. 
We can write the path as two coordinate functions $\gamma(t) = (x(t), y(t))$. Since $\gamma$ is continuous, $x(t)$ and $y(t)$ must be continuous.

By the Intermediate Value Theorem, since $x(0) = 0$ and $x(1) = 1$, the function $x(t)$ must hit every x-coordinate in between. Let's look at the peaks of the sine wave, which occur at $x$-coordinates $\frac{1}{2\pi n + \pi/2}$. 
We can construct a sequence of times $(a_n)_{n=0}^\infty$ in $[0,1]$ such that:
\[ x(a_n) = \frac{1}{2\pi n + \pi/2} \]
Notice that the y-coordinate at all these peaks is exactly $y(a_n) = 1$. 
As $n \to \infty$, the x-coordinate goes to 0, so $a_n \to 0$. But look at the limits:
\[ \lim_{n \to \infty} \gamma(a_n) = (0, 1) \]
However, sequential continuity demands that $\lim \gamma(a_n) = \gamma(0) = (0,0)$. 
$(0,1) \ne (0,0)$, which is a massive contradiction! Thus, no such continuous path can exist.
\end{solution}

\hrulefill

\section{Inner Product and Norm}

\begin{definition*}[Inner Product]
Let $V$ be a real vector space. An \textbf{inner product} on $V$ is a map $\langle \cdot, \cdot \rangle: V \times V \to \R$ satisfying:
\begin{enumerate}
    \item \textbf{Bilinearity:} $\langle \alpha u + \beta v, w \rangle = \alpha \langle u, w \rangle + \beta \langle v, w \rangle$ \textcolor{eGray}{(and equivalently in the second variable)}
    \item \textbf{Symmetry:} $\langle u, v \rangle = \langle v, u \rangle$
    \item \textbf{Definiteness:} $\langle u, u \rangle \ge 0$, and $\langle u, u \rangle = 0 \iff u = 0$
\end{enumerate}
\end{definition*}

\begin{definition*}[Norm]
A \textbf{norm} on $V$ is a map $\norm{\cdot}: V \to [0, \infty)$ satisfying:
\begin{enumerate}
    \item \textbf{Definiteness:} $\norm{v} \ge 0$, and $\norm{v} = 0 \iff v = 0$
    \item \textbf{Homogeneity:} $\norm{\alpha v} = \abs{\alpha} \norm{v}$
    \item \textbf{Triangle Inequality:} $\norm{u + v} \le \norm{u} + \norm{v}$
\end{enumerate}
\end{definition*}

\subsection*{Geometric Meaning and Projections}
Every inner product naturally defines a norm via $\norm{v} := \sqrt{\langle v, v \rangle}$. 
In $\R^2$, the standard inner product geometrically yields the angle $\theta$ between two vectors:
\[ \langle x, y \rangle = \norm{x} \norm{y} \cos\theta \]
In a general inner product space, this formula actually serves as the \emph{definition} of an angle! 

This allows us to define the orthogonal projection of a vector $u$ onto a vector $v$:
\[ \pi_v(u) = \frac{\langle u, v \rangle}{\norm{v}^2} v \]
In words, $\pi_v(u)$ is the "shadow" of $u$ pointing exactly along the direction of $v$.

\begin{center}
\begin{tikzpicture}[scale=1.5, very thick]
    % Axes
    \draw[->, eBlack, thick] (0,0) -- (5,0);
    \draw[->, eBlack, thick] (0,0) -- (0,3);
    
    % Vectors
    \coordinate (V) at (4.5, 1.5);
    \coordinate (U) at (2.5, 2.8);
    \coordinate (Proj) at (3.3, 1.1); % Approx projection
    
    \draw[->, eSaddleBrown] (0,0) -- node[below right] {$v$} (V);
    \draw[->, eMediumBlue] (0,0) -- node[above left] {$u$} (U);
    
    % Projection elements
    \draw[->, eMediumOrchid] (0,0) -- node[below right, yshift=-0.1cm] {$\pi_v(u)$} (Proj);
    \draw[dashed, eMediumOrchid] (U) -- (Proj);
    
    % Angle
    \draw[eDarkGreen] (1,0.33) arc (18.4:48.2:1.05cm);
    \node[eDarkGreen] at (1.1, 0.7) {$\theta$};
    
    % Right angle square
    \draw[eBlack, thin] (3.15, 1.35) -- (3.35, 1.45) -- (3.45, 1.25);
    \fill[eBlack] (3.3, 1.35) circle (0.5pt);
\end{tikzpicture}
\end{center}

\subsection*{The Cauchy-Schwarz Inequality}
If we take the norm of the projection, we get:
\[ \norm{\pi_v(u)} = \norm{u} \abs{\cos\theta} = \frac{\abs{\langle u, v \rangle}}{\norm{v}} \]
Because the "shadow" of a vector can never be longer than the vector itself ($\abs{\cos\theta} \le 1$), it must hold that $\norm{\pi_v(u)} \le \norm{u}$. Moving the denominator over yields the legendary Cauchy-Schwarz Inequality:
\[ \abs{\langle u, v \rangle} \le \norm{u} \norm{v} \]

\gemininote{Hint for a rigorous algebraic proof: Use the definiteness property! We know that the length of the \qt{perpendicular part} of the vector must be $\ge 0$. So, expand the inner product of the perpendicular part with itself: $\langle u - \pi_v(u), u - \pi_v(u) \rangle \ge 0$. Expanding this term will elegantly spit out the Cauchy-Schwarz inequality!}


% Week 4
\setcounter{section}{3} % Set to appropriate section number for Week 4

\section{The Differential}

To understand how derivatives work in multiple dimensions, let's first recall the 1-dimensional case from Analysis 1. For a function $f: \R \to \R$ that is differentiable at $x_0 \in \R$, we can write its linear approximation as:
\[ f(x_0 + h) = f(x_0) + f'(x_0)h + o(h) \]
This is perfectly equivalent to the classic limit definition:
\[ \lim_{h \to 0} \frac{f(x_0 + h) - f(x_0)}{h} = f'(x_0) \quad \iff \quad \lim_{h \to 0} \frac{\abs{f(x_0 + h) - f(x_0) - f'(x_0)h}}{\abs{h}} = 0 \]

Now, let's upgrade to a multivariable function $F: \R^n \to \R^m$. We can no longer just multiply by a scalar slope $f'(x_0)$. Instead, the derivative must be a \textbf{linear map} that transforms an $n$-dimensional step $h$ into an $m$-dimensional change in output!

\begin{definition*}[Differentiability]
We say that $F: \R^n \to \R^m$ (or $F: U \to \R^m$ for an open subset $U \subseteq \R^n$) is \textbf{differentiable} at $x_0$ if there exists a linear map $DF_{x_0}: \R^n \to \R^m$ such that:
\[ F(x_0 + h) = F(x_0) + DF_{x_0}(h) + o(\abs{h}) \]
Or, written as a limit:
\[ \lim_{h \to 0} \frac{\abs{F(x_0 + h) - F(x_0) - DF_{x_0}(h)}}{\abs{h}} = 0 \]
\end{definition*}

\gemininote{Notice the absolute value bars in the denominator! We divide by the \emph{length} of the vector $h$. If you look closely at the term $\frac{F(x_0 + h) - F(x_0)}{\abs{h}}$, you can see that as $h \to 0$, this evaluates the rate of change along a specific direction $\frac{h}{\abs{h}}$. This is intimately connected to the directional derivative!}

\subsection*{The Jacobian Matrix}

Since $DF_{x_0}: \R^n \to \R^m$ is a linear map, it can be represented by an $m \times n$ matrix—provided we fix a basis. In the standard basis of $\R^n$, this matrix is the famous \textbf{Jacobian matrix}:
\[ DF_x \approx JF(x) = \begin{pmatrix} 
\frac{\partial F_1}{\partial x_1} & \dots & \frac{\partial F_1}{\partial x_n} \\ 
\vdots & \ddots & \vdots \\ 
\frac{\partial F_m}{\partial x_1} & \dots & \frac{\partial F_m}{\partial x_n} 
\end{pmatrix} \]

Here, $F(x) = (F_1(x), \dots, F_m(x))^T$, and the $i$-th partial derivative is defined as taking the 1D derivative along the $i$-th standard basis vector $e_i$:
\[ \partial_i F(x) = \lim_{s \to 0} \frac{F(x + se_i) - F(x)}{s} \]

If $F$ is fully differentiable, the partial derivatives are simply the columns of the Jacobian, and we can compute them by applying the differential to the basis vectors:
\[ \partial_i F(x) = DF_x(e_i) \]

\geminiwarning{Crucial Trap: Partial derivatives can exist even if $F$ is NOT differentiable! Partial derivatives only check the grid lines (x-axis, y-axis). A function could be perfectly smooth along the axes but have a massive tear or crease diagonally. Full differentiability requires the linear approximation $DF_x$ to work from \emph{all} approach directions simultaneously.}

\subsection*{Examples of Differentials}

\begin{example*}[1. A Curve in 2D Space]
Consider a path $\gamma: \R \to \R^2, t \mapsto \binom{\cos t}{\sin t}$. \\
Let's find the differential at $t = \frac{\pi}{4}$. For a curve, the Jacobian is just a column vector (the velocity vector):
\[ \gamma'\left(\frac{\pi}{4}\right) := D\gamma_{\pi/4} = \begin{pmatrix} \partial_t \gamma_1(\pi/4) \\ \partial_t \gamma_2(\pi/4) \end{pmatrix} = \begin{pmatrix} -\sin(\pi/4) \\ \cos(\pi/4) \end{pmatrix} = \begin{pmatrix} -1/\sqrt{2} \\ 1/\sqrt{2} \end{pmatrix} \]
Using the approximation formula $F(x_0 + h) = F(x_0) + DF_{x_0}(h) + o(\abs{h})$:
\[ \gamma\left(\frac{\pi}{4} + h\right) = \frac{1}{\sqrt{2}}\binom{1}{1} + \frac{h}{\sqrt{2}}\binom{-1}{1} + o(\abs{h}) \]

\begin{center}
\begin{tikzpicture}[scale=2, very thick]
    % Axes
    \draw[->, eBlack] (-1.5, 0) -- (1.5, 0) node[right] {$x$};
    \draw[->, eBlack] (0, -1.5) -- (0, 1.5) node[above] {$y$};
    
    % Circle
    \draw[eMediumBlue] (0,0) circle (1cm);
    
    % Point at pi/4
    \coordinate (P) at (0.707, 0.707);
    \fill[eSaddleBrown] (P) circle (1.5pt) node[right, yshift=2mm] {$\gamma(\pi/4)$};
    
    % Tangent vector
    \draw[->, eDarkGreen] (P) -- ++(-0.707, 0.707) node[above] {$D\gamma_{\pi/4} = \gamma'(\pi/4)$};
\end{tikzpicture}
\end{center}
\textcolor{eGray}{Notice how the column vectors of the Jacobian are the partial derivatives. If we have a general direction vector $v = (v_1, \dots, v_n)^T$, the directional derivative is just the matrix-vector product: $\partial_v F(x) = DF_x(v) = JF(x)v = \sum_{i=1}^n v_i \frac{\partial F}{\partial x_i}$.}
\end{example*}

\begin{example*}[2. A Scalar Field]
Consider $f: \R^3 \to \R$, given by $(x,y,z) \mapsto 2x^2 + y^2 + 3z^2 - 2xyz$. \\
What is the differential $df(x,y,z)$? Since the output is 1D, the Jacobian is a row vector:
\[ Jf(x,y,z) = \left( 4x - 2yz, \quad 2y - 2xz, \quad 6z - 2xy \right) \]
Notice that this is exactly the transpose of the gradient!
\[ Jf(x,y,z) = \left( \frac{\partial f}{\partial x}, \frac{\partial f}{\partial y}, \frac{\partial f}{\partial z} \right) = (\nabla f)^T \]
\end{example*}

\hrulefill

\section{The Chain Rule}

For differentiable functions $f: U \subseteq \R^n \to \R^m$ and $g: V \subseteq \R^m \to \R^l$, the differential of the composition is the composition of the differentials:
\[ D(g \circ f)(x) = Dg(f(x)) \cdot Df(x) \]
Since linear maps compose via matrix multiplication, this translates directly to their Jacobians:
\[ J(g \circ f)(x) = Jg(f(x)) \cdot Jf(x) \]

\begin{exercise*}
Compute $J(g \circ f)$ for the following functions:
\begin{itemize}
    \item $f: \R^3 \to \R^2, \quad x \mapsto (x_1 x_2, \; x_2 + x_3)^T$
    \item $g: \R^2 \to \R^2, \quad y \mapsto (e^{y_1}, \; y_1 y_2)^T$
\end{itemize}
\end{exercise*}

\begin{solution}
First, compute the individual Jacobians:
\[ Jf(x) = \begin{pmatrix} x_2 & x_1 & 0 \\ 0 & 1 & 1 \end{pmatrix}, \quad Jg(y) = \begin{pmatrix} e^{y_1} & 0 \\ y_2 & y_1 \end{pmatrix} \]
Before multiplying, we must evaluate $Jg$ at the point $y = f(x)$, so $y_1 = x_1 x_2$ and $y_2 = x_2 + x_3$:
\[ Jg(f(x)) = \begin{pmatrix} e^{x_1 x_2} & 0 \\ x_2 + x_3 & x_1 x_2 \end{pmatrix} \]
Now, multiply the matrices:
\begin{equation}
    \begin{aligned}
        J(g \circ f)(x) &= \begin{pmatrix} e^{x_1 x_2} & 0 \\ x_2 + x_3 & x_1 x_2 \end{pmatrix} \begin{pmatrix} x_2 & x_1 & 0 \\ 0 & 1 & 1 \end{pmatrix} \\
        &= \begin{pmatrix} x_2 e^{x_1 x_2} & x_1 e^{x_1 x_2} & 0 \\ x_2(x_2 + x_3) & x_1(x_2 + x_3) + x_1 x_2 & x_1 x_2 \end{pmatrix}
    \end{aligned}
\end{equation}
\end{solution}

\subsection*{Alternative: Directional Derivatives}
Sometimes the differential is not easily expressed as a Jacobian matrix (e.g., when the vector space is the space of matrices). In such cases, it is extremely convenient to use the formula for directional derivatives:
\[ DF_{x_0}(v) = \partial_v F(x_0) = \frac{d}{ds}\bigg|_{s=0} F(x_0 + sv) \]
\textcolor{eGray}{Note: The directional derivative $\partial_v F(x_0)$ can exist even if $F$ is not fully differentiable at $x_0$. However, if $DF_{x_0}$ exists, this formula is always valid and often much faster to compute!}

\begin{exercise*}[Matrix Inner Products]
We identify the space of real-valued matrices $\R^{n \times n}$ with $\R^{n^2}$ using the standard inner product $\langle A, B \rangle = \text{Tr}(A^T B)$. 
Consider the matrix function:
\[ F: \R^{n \times n} \to \R^{n \times n}, \quad A \mapsto A^T A \]
Compute the differential at the identity matrix $Id \in \R^{n \times n}$, applied to a direction matrix $X \in \R^{n \times n}$:
\[ DF_{Id}(X) = \partial_X F(Id) \]
\end{exercise*}

\begin{solution}
Using the directional derivative trick, we evaluate the curve $Id + sX$ as it passes through $Id$:
\begin{equation}
    \begin{aligned}
        DF_{Id}(X) &= \frac{d}{ds}\bigg|_{s=0} F(Id + sX) \\
        &= \frac{d}{ds}\bigg|_{s=0} (Id + sX)^T (Id + sX) \\
        &= \frac{d}{ds}\bigg|_{s=0} [Id^T Id + sX^T Id + s Id^T X + s^2 X^T X] \\
        &= \frac{d}{ds}\bigg|_{s=0} [Id + s(X^T + X) + s^2 X^T X]
    \end{aligned}
\end{equation}
Taking the derivative with respect to $s$ and plugging in $s=0$ eliminates the constant term and the $s^2$ term, leaving exactly the linear part:
\[ DF_{Id}(X) = X^T + X \]
\end{solution}

\hrulefill

\section{Taylor Expansions in $\R^n$}

In the lecture, you saw the terrifying full formula for a function $f \in C^k(U, \R)$:
\[ f(\bar{x} + h) = \sum_{\abs{\alpha}=0}^k \partial^\alpha f(\bar{x}) \frac{h^\alpha}{\alpha!} + o(\abs{h}^k) \]
where $\alpha \in \N^n$ is a \qt{multi-index}.

\gemininote{What on earth is a multi-index? It's just a compact way to track all the crazy mixed partial derivatives! $\alpha = (\alpha_1, \dots, \alpha_n)$ tells you how many times to differentiate with respect to each variable. The sum $\abs{\alpha} = \alpha_1 + \dots + \alpha_n$ is the total order of the derivative. For example, $\partial^{(2,1)} = \partial_x^2 \partial_y$.}

\begin{exercise*}
Let $f \in C^3(\R^2, \R)$. Write out the maximal Taylor approximation (up to 3rd order) using multi-indices.
\end{exercise*}

\begin{solution}
Let's list all possible multi-indices $\alpha = (\alpha_1, \alpha_2)$ up to order 3:
\begin{itemize}
    \item $\abs{\alpha}=0$: $(0,0)$
    \item $\abs{\alpha}=1$: $(1,0), (0,1)$
    \item $\abs{\alpha}=2$: $(2,0), (1,1), (0,2)$
    \item $\abs{\alpha}=3$: $(3,0), (2,1), (1,2), (0,3)$
\end{itemize}
Recall the factorial rule: $\alpha! = \alpha_1! \alpha_2!$. For example, $(2,1)! = 2! \cdot 1! = 2$.
Plugging this into the formula yields an explosion of terms:
\begin{equation}
    \begin{aligned}
        f(\bar{x}+h) = \;&f(\bar{x}) + \partial_1 f(\bar{x})h_1 + \partial_2 f(\bar{x})h_2 \\
        &+ \frac{1}{2}\partial_1^2 f(\bar{x})h_1^2 + \partial_1 \partial_2 f(\bar{x})h_1 h_2 + \frac{1}{2}\partial_2^2 f(\bar{x})h_2^2 \\
        &+ \frac{1}{6}\partial_1^3 f(\bar{x})h_1^3 + \frac{1}{2}\partial_1^2 \partial_2 f(\bar{x})h_1^2 h_2 + \frac{1}{2}\partial_1 \partial_2^2 f(\bar{x})h_1 h_2^2 + \frac{1}{6}\partial_2^3 f(\bar{x})h_2^3 \\
        &+ o(\abs{h}^3)
    \end{aligned}
\end{equation}
\end{solution}

\subsection*{The Working Formula and Substitution Trick}
Because the multi-index formula is a nightmare to write out, we often rewrite it using the gradient operator as a polynomial:
\[ f(\bar{x} + h) = \sum_{l=0}^k \frac{1}{l!} \left( \sum_{m=1}^n h_m \partial_m \right)^l f(\bar{x}) + o(\abs{h}^k) \]

However, the absolute best way to compute Taylor series is often to \textbf{cheat} using known 1D series from Analysis 1!

\begin{exercise*}
Compute the Taylor expansion of $f(x,y) = (x + y^2)e^{-x^2-y^2}$ around $(0,0)$ up to third order.
\end{exercise*}

\begin{solution}
\textbf{Cumbersome solution:} Calculate all partial derivatives up to the 3rd order and plug them into the massive formula above. \textcolor{eDarkRed}{(Please never do this in an exam!)}

\textbf{The Easy Solution:} Use the substitution and multiplication rules from Analysis 1. We know the 1D Taylor series for $e^z = 1 + z + \frac{z^2}{2} + \dots$. Substitute $z = -x^2 - y^2$:
\[ e^{-x^2-y^2} = 1 + (-x^2 - y^2) + o(\abs{(x,y)}^3) \]
Now, multiply this by the front bracket, dropping any terms that result in a power higher than 3:
\begin{equation}
    \begin{aligned}
        f(x,y) &= (x + y^2) \left( 1 - x^2 - y^2 + o(\abs{(x,y)}^3) \right) \\
        &= x + y^2 - x^3 - xy^2 + o(\abs{(x,y)}^3)
    \end{aligned}
\end{equation}
And we are done in two lines of algebra!
\end{solution}

\vspace{0.5cm}
\begin{remark}
The 2nd order Taylor approximation of $f \in C^2(\R^n, \R)$ around $x_0 \in \R^n$ can be expressed beautifully using matrices:
\[ f(x_0 + h) = f(x_0) + \nabla f(x_0)^T h + \frac{1}{2} h^T \mathscr{H}f(x_0) h + o(\abs{h}^2) \]
where $\mathscr{H}f(x_0) = (\partial_i \partial_j f(x_0))_{i,j}$ is the symmetric \textbf{Hessian matrix} of $f$. 
\textcolor{eGray}{(Notice how perfectly this mirrors the 1D formula: $f(x_0) + f' h + \frac{1}{2}f'' h^2$!)}
\end{remark}


% Week 5
\setcounter{section}{4} % Setting counter for Week 5

\section{Optimization}

Let's look at finding the extreme values of a function $f: U \subseteq \R^n \to \R$. 

\begin{definition*}[Local Extrema and Critical Points]
\begin{itemize}
    \item $f$ has a \textbf{local minimum} at $x_0 \in U$ if there exists a ball $B_\eps(x_0)$ such that $\forall x \in B_\eps(x_0) \cap U$, we have $f(x) \ge f(x_0)$. \textcolor{eGray}{(Analogous for local maximum with $\le$).}
    \item Extrema are \textbf{strict} if the inequality is strict ($<$ or $>$).
    \item A \textbf{critical point} of $f \in C^1(U, \R)$ is any point $x_0$ where the gradient vanishes: $\nabla f(x_0) = 0$.
\end{itemize}
\end{definition*}

\begin{proposition*}
If $f \in C^1(U, \R)$ has a local extremum at an \emph{interior} point $x_0 \in \operatorname{int}(U)$, then $x_0$ must be a critical point ($\nabla f(x_0) = 0$).
\end{proposition*}

\gemininote{Geometrically, the gradient $\nabla f(x_0)$ always points in the direction of the \qt{steepest ascent}. If you are at a local minimum (the bottom of a valley) or a local maximum (the peak of a mountain), there is no \qt{uphill} direction anymore. You are standing perfectly flat. Therefore, the gradient vector must be the zero vector!}

\begin{center}
\begin{tikzpicture}[scale=1.5, very thick]
    % Coordinate grid
    \draw[eGray, dashed] (-1,-1) grid (3,2);
    \draw[->, eBlack] (-1.5, 0) -- (3.5, 0) node[right] {$x$};
    \draw[->, eBlack] (0, -1.5) -- (0, 2.5) node[above] {$y$};
    
    % Contour lines (Level sets)
    \draw[eMediumOrchid] (1,0.5) ellipse (1.5cm and 0.8cm);
    \draw[eMediumBlue] (1,0.5) ellipse (0.8cm and 0.4cm);
    
    % Minimum point
    \fill[eDarkGreen] (1,0.5) circle (2pt) node[below right] {Minimum};
    
    % Gradient vector
    \coordinate (P) at (2.2, 0.74); % Point on the outer contour
    \fill[eSaddleBrown] (P) circle (1.5pt);
    \draw[->, eDarkRed, line width=1.5pt] (P) -- ++(0.8, 0.6) node[above] {$\nabla f$};
    \node[eDarkRed, align=center] at (3.5, 1.6) {\small Steepest \\ \small ascent};
\end{tikzpicture}
\end{center}

\hrulefill

\section{Constrained Optimization (Lagrange)}

Suppose we want to find the hottest place on the surface of the Earth. 
In theory, temperature is a scalar field defined everywhere in space: $T \in C^\infty(\R^3, (0, \infty))$. We model the Earth's surface as the unit sphere $S^2 := \set{x \in \R^3 : \abs{x}^2 = 1}$.

We want to find the local maxima of the \emph{restricted} function $T|_{S^2}$. 
Crucially, these are \textbf{not necessarily} critical points of the global function $T$! 

\gemininote{Imagine standing on the equator. The true gradient of the temperature $\nabla T$ might point straight up into the sun (off the surface of the Earth). However, since you are constrained to walk only \emph{along} the surface, you only care about the \qt{shadow} of that gradient on the ground. A point is a constrained extremum if the true gradient points entirely off the surface, meaning it is perfectly perpendicular to the ground!}

Since our constraint surface $S^2$ is the level set of a function $g(x) = \abs{x}^2 - 1 = 0$, we know that $\nabla g(x)$ is everywhere perpendicular to the surface. 
Thus, at a constrained extremum $p$, the temperature gradient $\nabla T(p)$ must be perfectly parallel to the constraint gradient $\nabla g(p)$.

\begin{proposition*}[Lagrange Multipliers]
Suppose $f: U \subseteq \R^n \to \R$ and $g = (g_1, \dots, g_k): U \to \R^k$ are $C^1$. 
If the restricted function $f|_{g^{-1}\set{0}}$ has a local extremum at $x_0$, then the gradients $\set{\nabla f(x_0), \nabla g_1(x_0), \dots, \nabla g_k(x_0)}$ are linearly dependent.

Assuming the constraint gradients are independent, this means $\nabla f(x_0)$ is a linear combination of the constraint gradients. We can find these points by looking for critical points of the \textbf{Lagrangian}:
\[ L(x, \lambda) := f(x) - \sum_{i=1}^k \lambda_i g_i(x) \]
We force $\nabla L(x, \lambda) = 0$, which yields the system:
\begin{align*}
    \nabla f(x) - \lambda \cdot \nabla g(x) &= 0 \quad \text{\textcolor{eGray}{(Parallel gradients)}} \\
    g_j(x) &= 0 \quad \text{\textcolor{eGray}{(You are actually on the surface)}}
\end{align*}
\end{proposition*}

\begin{exercise*}
Find the local extrema of $f(x,y,z) = xyz$ on the sphere $K = \set{(x,y,z) \in \R^3 : x^2 + y^2 + z^2 \le 1}$.
\end{exercise*}

\begin{solution}
First, we check the \textbf{interior} ($x^2 + y^2 + z^2 < 1$) for standard critical points:
\[ \nabla f(x,y,z) = (yz, xz, xy) = 0 \implies \text{at least two variables must be } 0. \]
For example, $(0,0,0)$, or $(x,0,0)$. All these yield $f = 0$.

Next, we check the \textbf{boundary} using Lagrange multipliers. The constraint is $g(x,y,z) = x^2 + y^2 + z^2 - 1 = 0$. The Lagrangian is:
\[ L(x,y,z,\lambda) = xyz - \lambda(x^2 + y^2 + z^2 - 1) \]
Taking the partial derivatives and setting them to zero gives our system:
\begin{enumerate}
    \item $\partial_x L = yz - 2\lambda x = 0 \implies yz = 2\lambda x$
    \item $\partial_y L = xz - 2\lambda y = 0 \implies xz = 2\lambda y$
    \item $\partial_z L = xy - 2\lambda z = 0 \implies xy = 2\lambda z$
    \item $\partial_\lambda L = 1 - (x^2 + y^2 + z^2) = 0 \implies x^2 + y^2 + z^2 = 1$
\end{enumerate}

\gemininote{We want to divide by $x$, $y$, and $z$ to isolate $2\lambda$. Is that allowed? If $x=0$, equation (1) forces $y=0$ or $z=0$. If $x=0$ and $y=0$, equation (4) forces $z = \pm 1$. The points $(0,0,\pm 1)$ yield $f=0$, which we already know are not strict extrema since we can easily find positive/negative values elsewhere. Thus, we can safely assume $x, y, z \ne 0$ for the true extrema and divide!}

Assuming non-zero coordinates, we isolate $2\lambda$:
\[ 2\lambda = \frac{yz}{x} = \frac{xz}{y} = \frac{xy}{z} \]
Multiplying the terms pairwise yields:
\[ y^2 z^2 = x^2 z^2 = x^2 y^2 \implies x^2 = y^2 = z^2 \]
Plugging this symmetry into the constraint (4):
\[ 3x^2 = 1 \implies x = \pm \frac{1}{\sqrt{3}} \]
This gives $2^3 = 8$ solutions on the sphere:
\[ (x,y,z) = \left( \frac{\alpha}{\sqrt{3}}, \frac{\beta}{\sqrt{3}}, \frac{\gamma}{\sqrt{3}} \right) \quad \text{where } \alpha, \beta, \gamma \in \set{-1, 1} \]
\end{solution}

\hrulefill

\section{The Hessian Test}

Just like the second derivative test in Analysis 1, we can classify critical points in $\R^n$ using the second derivatives. The \textbf{Hessian matrix} of $f \in C^2(U, \R)$ is:
\[ \mathscr{H}f(x) = (\partial_i \partial_j f)_{i,j=1,\dots,n} \]
By Schwarz's Lemma, this matrix is perfectly symmetric.

\begin{definition*}[Definiteness of a Symmetric Matrix]
Let $A \in \R^{n \times n}$ be symmetric. It is called:
\begin{itemize}
    \item \textbf{Positive definite ($A > 0$):} $v^T A v > 0 \quad \forall v \ne 0$
    \item \textbf{Negative definite ($A < 0$):} $v^T A v < 0 \quad \forall v \ne 0$
    \item \textbf{Indefinite:} $\exists v, u$ such that $v^T A v > 0$ and $u^T A u < 0$
\end{itemize}
\end{definition*}

If $x_0$ is a critical point ($\nabla f(x_0) = 0$), then:
\begin{itemize}
    \item $\mathscr{H}f(x_0)$ is positive definite $\implies$ \textcolor{eDarkGreen}{\textbf{Local Minimum}}
    \item $\mathscr{H}f(x_0)$ is negative definite $\implies$ \textcolor{eDarkRed}{\textbf{Local Maximum}}
    \item $\mathscr{H}f(x_0)$ is indefinite $\implies$ \textcolor{eSaddleBrown}{\textbf{Saddle Point}}
\end{itemize}

\subsection*{How to determine Definiteness?}
Let $\lambda_1, \dots, \lambda_n$ be the eigenvalues of $A$. Recall that $\det(A) = \prod \lambda_i$ and $\operatorname{Tr}(A) = \sum \lambda_i$.
\begin{itemize}
    \item \textbf{For $n=2$:} 
        \begin{itemize}
            \item $\det(A) < 0 \implies$ Indefinite
            \item $\det(A) > 0$ and $\operatorname{Tr}(A) > 0 \implies$ Positive definite
            \item $\det(A) > 0$ and $\operatorname{Tr}(A) < 0 \implies$ Negative definite
        \end{itemize}
    \item \textbf{For $n \ge 3$ (The Hurwitz Criterion):} Check the leading principal minors (the determinants of the upper-left submatrices $A_1, A_2, \dots, A_n$).
        \begin{itemize}
            \item $A_1 > 0, A_2 > 0, A_3 > 0 \implies$ Positive definite
            \item $A_1 < 0, A_2 > 0, A_3 < 0 \implies$ Negative definite (alternating signs!)
            \item Anything else (excluding zeros) $\implies$ Indefinite
        \end{itemize}
\end{itemize}

\begin{example*}[Testing Matrices]
\textbf{1.} $A = \begin{pmatrix} 2 & 2 \\ 2 & 9 \end{pmatrix}$. $\det = 18 - 4 = 14 > 0$. $\operatorname{Tr} = 11 > 0 \implies$ \textcolor{eDarkGreen}{Positive definite}.

\textbf{2.} $A = \begin{pmatrix} 4 & 2 \\ 2 & -4 \end{pmatrix}$. $\det = -16 - 4 = -20 < 0 \implies$ \textcolor{eSaddleBrown}{Indefinite}.

\geminiwarning{The handwritten notes contained a matrix that was not symmetric ($\begin{smallmatrix} -5 & 0 & 0 \\ 0 & -4 & -7 \\ 3 & -2 & -4 \end{smallmatrix}$). The Hurwitz criterion strictly applies \emph{only} to symmetric matrices! Let's look at the correct symmetric matrix the TA meant to evaluate:}

\textbf{3.} $A = \begin{pmatrix} -5 & 0 & 0 \\ 0 & -4 & -2 \\ 0 & -2 & -4 \end{pmatrix}$. Let's check the Hurwitz minors:
\begin{align*}
    A_1 &= -5 \quad \textcolor{eGray}{(< 0)} \\
    A_2 &= \det\begin{pmatrix} -5 & 0 \\ 0 & -4 \end{pmatrix} = 20 \quad \textcolor{eGray}{(> 0)} \\
    A_3 &= \det(A) = -5 \cdot (16 - 4) = -60 \quad \textcolor{eGray}{(< 0)}
\end{align*}
The signs alternate ($-, +, -$), so this is \textcolor{eDarkRed}{Negative definite}.

\textbf{4.} $A = \begin{pmatrix} -1 & 0 & -4 \\ 0 & 2 & -2 \\ -4 & -2 & 22 \end{pmatrix}$. 
$A_1 = -1$ and $A_2 = \det\begin{pmatrix} -1 & 0 \\ 0 & 2 \end{pmatrix} = -2$. Since the signs are $(-, -, \dots)$, it violates both the positive and negative definite patterns $\implies$ \textcolor{eSaddleBrown}{Indefinite}.
\end{example*}

\hrulefill

\begin{exercise*}
Find the critical points of $f_\alpha(x,y) = x^3 - y^3 + 3\alpha xy$ for a constant $\alpha \in \R \setminus \set{0}$, and classify them (min/max/saddle).
\end{exercise*}

\begin{solution}
First, find the critical points by setting the gradient to zero:
\begin{align}
    \partial_x f &= 3x^2 + 3\alpha y = 0 \implies y = -\frac{x^2}{\alpha} \\
    \partial_y f &= -3y^2 + 3\alpha x = 0
\end{align}
Substitute (1) into (2):
\[ -3\left(-\frac{x^2}{\alpha}\right)^2 + 3\alpha x = 0 \implies -3\frac{x^4}{\alpha^2} + 3\alpha x = 0 \]
Divide by 3 and factor out $x$:
\[ x \left( \alpha - \frac{x^3}{\alpha^2} \right) = 0 \]
This yields two cases:
\begin{itemize}
    \item $x = 0 \implies y = 0$. Point: \textbf{$(0,0)$}.
    \item $x^3 = \alpha^3 \implies x = \alpha \implies y = -\alpha$. Point: \textbf{$(\alpha, -\alpha)$}.
\end{itemize}

Now, compute the Hessian matrix:
\[ \mathscr{H}f(x,y) = \begin{pmatrix} 6x & 3\alpha \\ 3\alpha & -6y \end{pmatrix} \]

\textbf{Test point 1: $(0,0)$}
\[ \mathscr{H}f(0,0) = \begin{pmatrix} 0 & 3\alpha \\ 3\alpha & 0 \end{pmatrix} \]
$\det = -9\alpha^2 < 0 \implies$ \textcolor{eSaddleBrown}{\textbf{Saddle Point}}.

\textbf{Test point 2: $(\alpha, -\alpha)$}
\[ \mathscr{H}f(\alpha, -\alpha) = \begin{pmatrix} 6\alpha & 3\alpha \\ 3\alpha & 6\alpha \end{pmatrix} = 3\alpha \begin{pmatrix} 2 & 1 \\ 1 & 2 \end{pmatrix} \]
The determinant is $36\alpha^2 - 9\alpha^2 = 27\alpha^2 > 0$. The trace is $12\alpha$.
\begin{itemize}
    \item If $\alpha > 0$: Trace is positive $\implies$ Positive definite $\implies$ \textcolor{eDarkGreen}{\textbf{Local Minimum}}.
    \item If $\alpha < 0$: Trace is negative $\implies$ Negative definite $\implies$ \textcolor{eDarkRed}{\textbf{Local Maximum}}.
\end{itemize}
\end{solution}

% Week 6
\setcounter{section}{5} % Set to appropriate section number for Week 6

\section{Convexity}

\begin{definition*}[Convex Sets and Functions]
\begin{itemize}
    \item A subset $A \subseteq \R^n$ is \textbf{convex} if for any $x, y \in A$ and any $t \in [0,1]$, the point $p = (1-t)x + ty$ is also an element of $A$. \textcolor{eGray}{(The straight line segment connecting any two points in the set lies entirely within the set.)}
    
    \item A function $f: A \to \R$ is \textbf{convex} if its domain $A$ is a convex set, and for all $x, y \in A$ and $t \in [0,1]$:
    \[ f((1-t)x + ty) \le (1-t)f(x) + t f(y) \]
\end{itemize}
\end{definition*}

\gemininote{The functional inequality simply states that if you draw a straight secant line between two points on the graph of $f$, the actual graph of the function will always hang \emph{below} (or exactly on) that secant line. It looks like a bowl.}

\begin{proposition*}
Let $U \subseteq \R^n$ be open and convex, and $f \in C^2(U)$. The following statements are perfectly equivalent:
\begin{enumerate}
    \item $f$ is a convex function.
    \item The Hessian matrix $\mathscr{H}f(x)$ is positive semi-definite (non-negative definite) for all $x \in U$.
    \item $f(y) - f(x) \ge Df_x(y-x)$ for all $x, y \in U$. \textcolor{eGray}{(The graph of the function always lies above any of its tangent planes.)}
\end{enumerate}
\end{proposition*}

\begin{example*}
The $n$-dimensional parabola $p: \R^n \to \R, x \mapsto \abs{x}^2 = \sum_{i=1}^n x_i^2$ is convex. 
Taking the second derivatives, we get $\partial_i \partial_j p(x) = 2\delta_{ij}$ (which is $2$ if $i=j$, and $0$ otherwise). Thus, the Hessian is strictly diagonal:
\[ \mathscr{H}p(x) = 2 I_n \]
Every eigenvalue of this matrix is $2 > 0$, so the matrix is strictly positive definite. This immediately implies the function is convex!
\end{example*}

\hrulefill

\begin{exercise*}
Let $U \subseteq \R^n$ be open and convex, and $f \in C^1(U)$.
\begin{enumerate}
    \item Show that $f: U \to \R$ is convex if and only if the 1-dimensional slice functions $f_{x,y}: [0,1] \to \R, s \mapsto f(x + s(y-x))$ are convex for all $x, y \in U$.
    \item Use this to prove the tangent plane property: $f$ convex $\iff f(y) - f(x) \ge Df_x(y-x)$.
    \item Conclude that if a convex function $f$ has a critical point at $x_0 \in U$, then it has a \emph{global} minimum at $x_0$.
\end{enumerate}
\end{exercise*}

\begin{solution}\hfill\\
\textbf{a) The 1D Reduction} \\
Let $f_{x,y}$ be convex. We want to show $f$ is convex.
\begin{align*}
    f_{x,y}((1-t)s_0 + t s_1) &\le (1-t)f_{x,y}(s_0) + t f_{x,y}(s_1) \\
    f(x + [(1-t)s_0 + t s_1](y-x)) &\le (1-t)f(x + s_0(y-x)) + t f(x + s_1(y-x)) 
\end{align*}
If we simply set $s_0 = 0$ and $s_1 = 1$, the right side simplifies to $(1-t)f(x) + t f(y)$, and the left side simplifies exactly to $f((1-t)x + ty)$. This immediately proves $f$ is convex!
Vice versa, if $f$ is convex, the convexity of $f_{x,y}$ follows by applying the convexity of $f$ to the two points $u = x + s_0(y-x)$ and $v = x + s_1(y-x)$.

\vspace{0.3cm}
\textbf{b) Tangent Plane Inequality} \\
Recall from Analysis 1 that a 1D function $g \in C^1(I)$ is convex if and only if its derivative $g'$ is monotonically increasing. 
By part (a), $f$ is convex $\iff f_{x,y}$ is convex $\iff f_{x,y}'$ is increasing $\forall x,y \in U$.

Compute the derivative of the slice using the chain rule for $t \in [0,1]$:
\[ f_{x,y}'(t) = Df_{x + t(y-x)}(y-x) \]
By the 1D Mean Value Theorem, there must exist some $s \in (0,1)$ such that:
\[ f(y) - f(x) = f_{x,y}(1) - f_{x,y}(0) = f_{x,y}'(s) = Df_{x + s(y-x)}(y-x) \]
Now we apply the increasing property. If $f$ is convex, then $f_{x,y}'(s) \ge f_{x,y}'(0)$.
\[ f(y) - f(x) = f_{x,y}'(s) \ge f_{x,y}'(0) = Df_x(y-x) \]
\textcolor{eGray}{(The proof for the reverse direction involves evaluating the inequality for intermediate points $u$ and $v$ to show that $f_{x,y}'$ must be increasing. See the handwritten notes for the algebraic manipulation.)}

\vspace{0.3cm}
\textbf{c) Critical Points of Convex Functions} \\
Since $f$ is convex, we can always apply the inequality we just proved. If $x_0$ is a critical point, its differential is zero ($Df_{x_0} = 0$). Plugging this into the inequality for any arbitrary point $y \in U$:
\[ f(y) - f(x_0) \ge Df_{x_0}(y - x_0) = 0 \implies f(y) \ge f(x_0) \]
Since this holds for \emph{all} $y \in U$, $x_0$ is an absolute global minimum!
\end{solution}

\hrulefill

\section{The Inverse Function Theorem (IFT)}

\begin{theorem*}[Inverse Function Theorem]
Let $U \subseteq \R^n$ be open and $f \in C^k(U, \R^n)$. Let $x_0 \in U$ be a point such that the differential $Df_{x_0}$ is strictly invertible (i.e., $\det(Jf(x_0)) \ne 0$).

Then there exists an open neighborhood $U_0 \subseteq U$ around $x_0$ such that:
\begin{enumerate}
    \item $f|_{U_0}: U_0 \to f(U_0)$ is a $C^k$-diffeomorphism. \textcolor{eGray}{(It is locally invertible and smooth!)}
    \item The differential of the inverse is the inverse of the differential: 
    \[ D(f|_{U_0}^{-1})_{f(x)} = \left( D(f|_{U_0})_x \right)^{-1} \quad \forall x \in U_0 \]
\end{enumerate}
\end{theorem*}

\gemininote{Intuition: If the derivative matrix is invertible, it means the function locally acts like a linear transformation that scales or rotates space, but it never \qt{flattens} a dimension down to zero. Because no information is destroyed (no dimensions are crushed), you can easily reverse the process—at least locally!}

\begin{question*}
\begin{enumerate}
    \item If $f: \R^n \to \R^n$ is $C^k$ and a homeomorphism (bijective with continuous inverse), is it necessarily a $C^k$-diffeomorphism?
    \item If $f \in C^\infty(U, \R^n)$ has an everywhere invertible differential ($Df_x$ invertible $\forall x \in U$), is $f$ a $C^\infty$-diffeomorphism onto its image?
\end{enumerate}
\end{question*}

\begin{answer*}
\textbf{1. False.} Consider $f: \R \to \R, x \mapsto x^3$. This function is bijective and continuous, and its inverse $f^{-1}(x) = x^{1/3}$ is also continuous. However, $f^{-1}$ is not differentiable at $x=0$, so it is not a diffeomorphism!

\textbf{2. False.} Local invertibility does not guarantee global invertibility! Consider $f(x) = \cos(x)$ restricted to $\R \setminus \pi\Z$. For any $x$ in this domain, $f'(x) = -\sin(x) \ne 0$, so the differential is everywhere invertible. The IFT guarantees it is locally a diffeomorphism everywhere. But globally, cosine oscillates and hits the same values infinitely many times, so it is obviously not bijective onto $(-1, 1)$!
\end{answer*}

\hrulefill

\section{The Implicit Function Theorem}

\begin{theorem*}[Implicit Function Theorem]
\textbf{Setup:} Let $n > n-d \ge 1$. Let $f \in C^k(\R^d \times \R^{n-d}, \R^{n-d})$. Suppose $f(x_0, y_0) = 0$ for some point. 
Suppose $D_y f(x_0, y_0)$, the differential with respect to the $y$-variables, is bijective (i.e., $J_y f(x_0, y_0)$ is an invertible square matrix).

\textbf{Consequence:} There exist open balls $B_r(x_0) \subseteq \R^d$ and $B_s(y_0) \subseteq \R^{n-d}$, and a $C^k$ function $g: B_r(x_0) \to B_s(y_0)$ such that for all points in this local window:
\[ f(x,y) = 0 \iff y = g(x) \]
Furthermore, the Jacobian of this new implicit function is:
\[ Jg(x) = -\left( J_y f(x, g(x)) \right)^{-1} J_x f(x, g(x)) \]
\end{theorem*}

\gemininote{The TA wrote: \qt{The geometrical meaning of this formula eludes me, sorry.} But it's actually incredibly logical! It's just the Multivariable Chain Rule. We know that locally, $f(x, g(x)) = 0$. If we take the total derivative of both sides with respect to $x$, the Chain Rule gives:
$J_x f + J_y f \cdot Jg = 0$. 
Now just solve for $Jg$! Move $J_x f$ to the other side, and multiply from the left by the inverse of $J_y f$. That's it!}

\begin{example*}[1. The Unit Circle]
Consider $F(x,y) = x^2 + y^2 - 1 = 0$. Can we express $y$ as a function of $x$?
Taking the $y$-derivative: $D_y F(x,y) = 2y$.
This is invertible (non-zero) everywhere \emph{except} at $y=0$ (the points $(-1,0)$ and $(1,0)$). 
Thus, for any point where $y \ne 0$, the IFT guarantees we can locally write $y(x) = \pm\sqrt{1-x^2}$. At $y=0$, the tangent is perfectly vertical, so it fails the vertical line test and cannot be written as $y(x)$.
\end{example*}

\begin{example*}[2. The Quintic Equation]
Consider the level set $G^{-1}\set{0} = \set{(x,y) \in \R^2 : y^5 + y^3 + y + x = 0}$.
Clearly, the origin $(0,0)$ is a solution since $0^5 + 0^3 + 0 + 0 = 0$.
Check the partial derivatives at $(0,0)$:
\begin{align*}
    \partial_x G(0,0) &= 1 \ne 0 \\
    \partial_y G(0,0) &= \left[ 5y^4 + 3y^2 + 1 \right]_{y=0} = 1 \ne 0
\end{align*}
Since $\partial_y G \ne 0$, the IFT guarantees that there exists an open interval around $x=0$ where we can uniquely write $y = g(x)$. 

\textcolor{eSaddleBrown}{Why is this theorem so powerful?} By the famous Abel-Ruffini Theorem from Galois theory, there is mathematically \textbf{no general algebraic formula} to solve a polynomial of degree 5. You literally cannot write down $y(x)$ using standard roots. Yet, the Implicit Function Theorem mathematically guarantees that this function $y(x)$ exists, that it is perfectly smooth, and we can even compute its derivative without ever knowing the function's explicit formula!
\end{example*}




% Week 7 - 

\setcounter{section}{6} 

\section{Submanifolds}

\textcolor{eGray}{Recall: A $C^k$-diffeomorphism is a bijective function whose inverse is also $C^k$. It essentially acts as a perfectly smooth transformation that stretches and bends space without ever tearing it or creating sharp creases.}

\begin{definition*}
$M^m \subseteq \R^n$ is a $C^k$, $m$-dimensional \emph{submanifold} of $\R^n$ if for any point $p \in M$ there are:
\begin{itemize}
    \item an \emph{open} neighborhood $U \subseteq \R^n$ of $p$
    \item an \emph{open} neighborhood $V \subseteq \R^n$ of $0$
    \item a $C^k$-diffeomorphism $\Phi: U \to V$ with $\Phi(p) = 0$ \textcolor{eGray}{(the submanifold chart)}
\end{itemize}
such that
\[ \Phi(U \cap M) = (\R^m \times \set{0}) \cap V = \set{x \in V : x_{m+1} = \dots = x_n = 0} \]
\end{definition*}

\gemininote{The intuition here is beautifully simple: No matter how curved or twisted $M$ looks in the big space $\R^n$, if you zoom in close enough to a point $p$ (inside the window $U$), you can find a magical map $\Phi$ that completely \qt{irons out} the wrinkles. It flattens the manifold $M$ so it perfectly aligns with the standard flat coordinate axes of $\R^m$.}

\begin{center}
\begin{tikzpicture}[scale=1.5, very thick]
    % U and M
    \draw[eDarkGray, dashed] (-1, 0) ellipse (1.5 and 1);
    \node[eDarkGray] at (-1, 1.2) {$U$};
    \draw[eMediumBlue] (-2.2, -0.2) to[out=20, in=160] (0.2, 0.2);
    \node[eMediumBlue] at (-1.8, -0.4) {$M$};
    \fill[eSaddleBrown] (-1, 0.05) circle (1.5pt) node[below] {$p$};

    % Phi Arrow
    \draw[->, eDarkRed] (0.8, 0) -- node[above] {$\Phi$} (2.2, 0);

    % V and R^m
    \draw[eDarkGray, dashed] (4, 0) ellipse (1.5 and 1);
    \node[eDarkGray] at (4, 1.2) {$V$};
    \draw[eMediumBlue] (2.2, 0) -- (5.8, 0);
    \node[eMediumBlue] at (5.2, -0.2) {$\R^m$};
    \fill[eSaddleBrown] (4, 0) circle (1.5pt) node[below] {$0$};
    \draw[->, eBlack] (4, -1) -- (4, 1) node[right] {$\R^{n-m}$};
\end{tikzpicture}
\end{center}

Finding such a chart $\Phi$ manually can be extremely difficult. It is often much easier to prove a set is a submanifold by using one of the two standard perspectives: \emph{Implicit Representation} (defining the manifold by what it restricts) or \emph{Local Parametrization} (defining the manifold by building it).

\subsection*{The Implicit Perspective: Cutting down space}

To separate the implicit definition from the parametric one, we will use \textcolor{eDarkGreen}{\textbf{green for the implicit map $F$}} and \textcolor{eMediumBlue}{\textbf{blue for the parametric map $f$}}.

\begin{theorem*}[Regular value theorem]
Suppose $U \subseteq \R^n$ is \emph{open}, $\textcolor{eDarkGreen}{F}: U \to \R^l$ is $C^k$, and $y \in \textcolor{eDarkGreen}{F}(U)$. If the differential
\[ D\textcolor{eDarkGreen}{F}_x: \R^n \to \R^l \]
is surjective (i.e. $J\textcolor{eDarkGreen}{F}(x)$ has full rank) for all $x \in \textcolor{eDarkGreen}{F}^{-1}\set{y} = \set{\bar{x} \in U : \textcolor{eDarkGreen}{F}(\bar{x}) = y}$, then the level set $\textcolor{eDarkGreen}{F}^{-1}\set{y}$ is a $C^k$ submanifold with dimension $n-l$.
\end{theorem*}

\gemininote{Why must it be surjective? Surjectivity of the derivative means the constraints defined by $\textcolor{eDarkGreen}{F}$ are truly independent and don't collapse on each other. If the rank dropped, the manifold might have sharp corners, self-intersections, or singular points where it's no longer \qt{smooth}.}

\begin{example*}
Consider $\textcolor{eDarkGreen}{F}: \R^2 \to \R$, given by $(x,y) \mapsto x^2 + y^2$. \\
Then the Jacobian is $J\textcolor{eDarkGreen}{F}(x,y) = (2x, 2y)$. \\
This matrix (a row vector here) has full rank (rank 1) as long as it isn't the zero vector. So it has full rank if $x \ne 0$ or $y \ne 0$, i.e., $\forall (x,y) \in \R^2 \setminus \set{0}$.

For any $R > 0$, the origin $(0,0) \notin \textcolor{eDarkGreen}{F}^{-1}\set{R^2}$. Therefore, the level set
\[ \textcolor{eDarkGreen}{F}^{-1}\set{R^2} = \set{(x,y) \in \R^2 : x^2 + y^2 = R^2} \]
is guaranteed to be a smooth submanifold of dimension $2 - 1 = 1$.

\begin{center}
\begin{tikzpicture}[scale=1.5, very thick]
    \draw[->, eBlack] (-1.5, 0) -- (1.5, 0) node[right] {$x$};
    \draw[->, eBlack] (0, -1.5) -- (0, 1.5) node[above] {$y$};
    \draw[eMediumOrchid] (0,0) circle (1cm);
    \draw[eDarkRed, -] (0,0) -- node[below right] {$R$} (45:1cm);
\end{tikzpicture}
\end{center}

Similarly, $\set{x \in \R^n : \abs{x}^2 = R^2}$ is a submanifold of dimension $n-1$ for $R > 0$. \\
\textcolor{eGray}{This is the classical $n$-sphere of radius $R$.} 
\geminiwarning{Kurzer Terminologie-Check: Eine Sphäre, die im \(\R^n\) lebt, ist streng genommen eine \((n-1)\)-dimensionale Fläche. Topologen nennen sie daher meist die \((n-1)\)-Sphäre, notiert als \(S^{n-1}\). Hier in den Notizen wird sie informell die \(n\)-Sphäre genannt (basierend auf der Dimension des Umgebungsraumes).}
\end{example*}

\subsection*{The Connection to Lagrangian Optimization}

Recall how we build the Lagrangian $L = f - \sum_{i=1}^k \lambda_i g_i$ when optimizing a function $f|_M$ for $f \in C^1(\R^n, \R)$. The constraint set $M$ is defined as:
\[ g_1, \dots, g_k: U \to \R, \quad M := \set{x \in \R^n : g_1(x) = \dots = g_k(x) = 0} \]
For the method of Lagrange multipliers to work, we require that the gradients $(\nabla g_1, \dots, \nabla g_k)$ are linearly independent $\forall x \in M$.

Let's package all these constraints into a single implicit map $\textcolor{eDarkGreen}{g}: \R^n \to \R^k$, where $x \mapsto (g_1(x), \dots, g_k(x))$. The requirement that the gradients are linearly independent means exactly that the Jacobian:
\[ J\textcolor{eDarkGreen}{g}(x) = \begin{pmatrix} \longleftarrow & \nabla g_1(x) & \longrightarrow \\ & \vdots & \\ \longleftarrow & \nabla g_k(x) & \longrightarrow \end{pmatrix} \]
has full rank! Therefore, the constraint surface $M = \textcolor{eDarkGreen}{g}^{-1}\set{0}$ is automatically an $(n-k)$-dimensional submanifold by the regular value theorem!

\gemininote{Die Notizen erwähnen hier \qt{ergo $k \le n$}. Why is this so crucial? Think about what happens if you add constraints. In a 3D room ($n=3$), one constraint ($k=1$) restricts you to a 2D floor. Two constraints ($k=2$) restrict you to a 1D line (where two floors intersect). Three constraints ($k=3$) restrict you to a 0D point. If you impose $k=4$ independent constraints in a 3D room, it's physically impossible—the set of solutions is empty! The rank of a matrix can never exceed its number of columns, so for $J\textcolor{eDarkGreen}{g}$ to have full row rank $k$, we mathematically must have $k \le n$.}


\subsection*{The Parametric Perspective: Building up space}

\begin{theorem*}[Local parametrization]
Let $M \subseteq \R^n$. If $\forall p \in M$ there exist an \emph{open} neighborhood $V \subseteq \R^n$ of $p$, an \emph{open} set $U \subseteq \R^m$ and a $C^k$ map $f: U \to V \cap M$ such that:
\begin{enumerate}
    \item $Df_x: \R^m \to \R^n$ is injective $\forall x \in U$ \textcolor{eGray}{(the map doesn't crush dimensions together)}
    \item $f$ is a homeomorphism \textcolor{eGray}{(bijective, continuous, and $f^{-1}$ is continuous, ensuring no self-intersections)}
\end{enumerate}
then $M$ is a $C^k$ submanifold with dimension $m$.
\end{theorem*}

\begin{example*}[Parabola]
Let $P = \set{(x,t) \in \R^n \times \R : t = \abs{x}^2}$ \textcolor{eGray}{(the $n$-parabola)}.

\begin{center}
\begin{tikzpicture}[scale=1.5, very thick]
    \draw[->, eBlack] (-1.5, 0) -- (1.5, 0) node[right] {$x$};
    \draw[->, eBlack] (0, -0.5) -- (0, 1.5) node[above] {$t$};
    \draw[eMediumBlue, domain=-1.2:1.2, samples=50] plot (\x, {\x*\x});
\end{tikzpicture}
\end{center}

We could define this implicitly and use the regular value theorem, but let's build it parametrically. Consider the map:
\[ f: \R^n \to \R^n \times \R, \quad x \mapsto (x, \abs{x}^2) \]
We check our conditions:
\begin{enumerate}
    \item $Jf(x) = \begin{pmatrix} Id_{n \times n} \\ 2x^\top \end{pmatrix}$ clearly has full column rank, so the derivative is injective.
    \item $f$ is bijective onto its image and continuous. Its inverse is the canonical projection $\pi: f(U) \to \R^n$, mapping $(x,t) \mapsto x$, which is trivially continuous.
    \item $f(\R^n) = P$.
\end{enumerate}
So $P$ is an $n$-dimensional submanifold.
\end{example*}

\textcolor{eGray}{Note: The graphical representation of a function (like plotting $y = h(x)$) is just a very common, equivalent special case of a local parametrization.}

\begin{example*}[Alternative optimization technique]
Let $f: \R^2 \to \R$, given by $(x,y) \mapsto 2x^2 + y^2 - x$. \\
Let's find the extrema of $f|_{S^1}$ \emph{without} using Lagrange multipliers!

Since $S^1$ is a 1-dimensional submanifold, we can just parametrize $S^1 = \set{(x,y) \in \R^2 : x^2 + y^2 = 1}$ using an angle $\theta$:
\[ \begin{pmatrix} x(\theta) \\ y(\theta) \end{pmatrix} = \begin{pmatrix} \cos\theta \\ \sin\theta \end{pmatrix} \]
We substitute this parametrization directly into our objective function:
\begin{equation}
    \begin{aligned}
        f(\theta) &= f(x(\theta), y(\theta)) = 2\cos^2\theta + \sin^2\theta - \cos\theta \\
        &= 2\cos^2\theta + (1 - \cos^2\theta) - \cos\theta \\
        &= \cos^2\theta - \cos\theta + 1
    \end{aligned}
\end{equation}
Now it's just a standard Analysis 1 problem! Setting the derivative to zero:
\[ f'(\theta) = -2\cos\theta\sin\theta + \sin\theta = 0 \]
Factoring out $\sin\theta$:
\[ \sin\theta(1 - 2\cos\theta) = 0 \]
This gives two cases where the product is zero:
\begin{enumerate}
    \item $\sin\theta = 0 \implies \theta_1 \in \set{0, \pi}$
    \item $(1 - 2\cos\theta) = 0 \implies \cos\theta = \frac{1}{2} \implies \theta_2 \in \set{\frac{\pi}{3}, \frac{5\pi}{3}}$
\end{enumerate}
Plugging these angles back into our parametrization $(x,y) = (\cos\theta, \sin\theta)$, the extrema on the circle are located at:
\begin{enumerate}
    \item $(\cos\theta_1, \sin\theta_1) = (\pm 1, 0)$
    \item $(\cos\theta_2, \sin\theta_2) = \left(\frac{1}{2}, \pm\frac{\sqrt{3}}{2}\right)$
\end{enumerate}
\end{example*}

\hrulefill

\section{Tangent and Normal space}

Suppose $M^m \subseteq \R^n$ is a submanifold and $f: U \to f(U) \subseteq M$ is a local parametrization around some point $p \in f(U)$ with $f(q) = p$.

The $n \times m$-matrix of partial derivatives
\[ Jf(q) = \begin{pmatrix} | & & | \\ \partial_1 f(q) & \dots & \partial_m f(q) \\ | & & | \end{pmatrix} \]
has full rank, meaning the column vectors $\partial_i f(q)$ are linearly independent. 

\gemininote{Think of the vectors $\partial_i f(q)$ as the \qt{velocity vectors} you get if you walk along the grid lines of your parameter space $U$. Since they are linearly independent, they span a perfectly flat $m$-dimensional plane that grazes the manifold at $p$.}

\begin{definition*}[Tangent space]
For $M, f$ as above, the Tangent space of $M$ at $p \in M$ is the space spanned by these velocity vectors:
\[ T_p M = \Span\set{\partial_1 f(q), \dots, \partial_m f(q)} = Df_q(\R^m) \]
Remarkably, this space is completely independent of the specific parametrization $f$ you chose to use!
\end{definition*}

\begin{example*}
We parametrize the upper hemisphere
\[ S_+^2 = \set{(x,y,z) \in \R^3 : x^2 + y^2 + z^2 = 1, z > 0} \]
graphically as a function of $x$ and $y$:
\[ f(x,y) = \begin{pmatrix} x \\ y \\ \sqrt{1 - x^2 - y^2} \end{pmatrix} \]
Let's determine the tangent space $T_{e_3} S_+^2$ at the \qt{North Pole}, where $e_3 = (0,0,1)$. This corresponds to the parameters $(x,y) = (0,0)$.

First, we calculate the partial derivatives and evaluate them at the origin:
\[ \partial_x f \Big|_{(x,y)=(0,0)} = \begin{pmatrix} 1 \\ 0 \\ \frac{-x}{\sqrt{1 - x^2 - y^2}} \end{pmatrix} \Bigg|_{(x,y)=(0,0)} = \begin{pmatrix} 1 \\ 0 \\ 0 \end{pmatrix} =: e_1 \]
\[ \partial_y f \Big|_{(x,y)=(0,0)} = \begin{pmatrix} 0 \\ 1 \\ \frac{-y}{\sqrt{1 - x^2 - y^2}} \end{pmatrix} \Bigg|_{(x,y)=(0,0)} = \begin{pmatrix} 0 \\ 1 \\ 0 \end{pmatrix} =: e_2 \]
So, the tangent space at the North Pole is perfectly flat against the xy-plane:
\[ T_{e_3} S_+^2 = T_{e_3} S^2 = \Span\set{e_1, e_2} = \R^2 \times \set{0} \]

\begin{center}
\begin{tikzpicture}[scale=2, very thick]
    % Sphere
    \draw[eMediumBlue] (0,0) circle (1cm);
    \draw[eMediumBlue, dashed] (1,0) arc (0:180:1cm and 0.3cm);
    \draw[eMediumBlue] (1,0) arc (0:-180:1cm and 0.3cm);

    % Tangent Plane
    \draw[eDarkGreen] (-1.2, 1) -- (0.5, 0.7) -- (1.2, 1.2) -- (-0.5, 1.5) -- cycle;
    \node[eDarkGreen] at (1.4, 1) {$T_{e_3} S^2$};

    % Point e_3
    \fill[eDarkGreen] (0,1) circle (1pt);
\end{tikzpicture}
\end{center}
\end{example*}

\subsection*{The Grand Synthesis: Tangent vs. Normal}

Now, let's look at the beautiful relationship between our two perspectives. Suppose $\textcolor{eDarkGreen}{F}: U^n \to \R^{n-m}$ is an \emph{implicit representation} of a submanifold $M^m = \textcolor{eDarkGreen}{F}^{-1}\set{0}$, and $f: V^m \to U$ is a \emph{local parametrization} of that same manifold.

Then, by definition of these two constructs, three things are true:
\begin{enumerate}
    \item $J\textcolor{eDarkGreen}{F}(x) = \begin{pmatrix} \longleftarrow & \nabla \textcolor{eDarkGreen}{F}_1(x) & \longrightarrow \\ & \vdots & \\ \longleftarrow & \nabla \textcolor{eDarkGreen}{F}_{n-m}(x) & \longrightarrow \end{pmatrix}$ has full rank. Thus, the gradients $\set{\nabla \textcolor{eDarkGreen}{F}_i(x)}_{i=1,\dots,n-m}$ are linearly independent.
    \item $T_{f(q)} M = \Span\set{\partial_i f(q)}$ as we just discussed.
    \item $\textcolor{eDarkGreen}{F}(f(x)) = 0 \quad \forall x \in V$. This is the crucial link: because the image of the parametrization $f$ lies entirely on the manifold, feeding $f(x)$ into the constraint equation $\textcolor{eDarkGreen}{F}$ must always output exactly zero!
\end{enumerate}

From point \textbf{c)}, we can apply the multivariable chain rule to see what happens to the derivatives:
\[ J(\textcolor{eDarkGreen}{F} \circ f)(q) = J\textcolor{eDarkGreen}{F}(f(q)) \cdot Jf(q) = \begin{pmatrix} \longleftarrow & \nabla \textcolor{eDarkGreen}{F}_1(f(q)) & \longrightarrow \\ & \vdots & \\ \longleftarrow & \nabla \textcolor{eDarkGreen}{F}_{n-m}(f(q)) & \longrightarrow \end{pmatrix} \begin{pmatrix} | & & | \\ \partial_1 f(q) & \dots & \partial_m f(q) \\ | & & | \end{pmatrix} = 0 \]

Now, look closely at what this matrix multiplication is actually doing \warning{(!!!)}:
\[ (J\textcolor{eDarkGreen}{F}(f(q)) \cdot Jf(q))_{ij} = \langle \nabla \textcolor{eDarkGreen}{F}_i(f(q)), \partial_j f(q) \rangle = 0 \quad \forall i=1,\dots,n-m, \quad j=1,\dots,m \]

Every single gradient vector of $\textcolor{eDarkGreen}{F}$ has a dot product of zero with every single tangent vector from $f$. So the vectors $\nabla \textcolor{eDarkGreen}{F}_i(f(q))$ are all strictly perpendicular (normal) to the tangent space $T_{f(q)} M$. 

Since by point \textbf{a)}, we have exactly $n-m$ linearly independent vectors, and they are all perpendicular to our $m$-dimensional tangent plane, they perfectly span the remaining dimensions of $\R^n$. This leads us to define:

\begin{definition*}[Normal space]
Let $M^m \subseteq \R^n$ be a submanifold and $\textcolor{eDarkGreen}{F}: U^n \to \R^{n-m}$ an implicit representation where $M = \textcolor{eDarkGreen}{F}^{-1}\set{0}$.
Then the normal space of $M$ at $p \in M$ is spanned by the gradients of the constraints:
\[ N_p M := \Span\set{\nabla \textcolor{eDarkGreen}{F}_1(p), \dots, \nabla \textcolor{eDarkGreen}{F}_{n-m}(p)} \]
\end{definition*}

\begin{example*}
Let's represent the sphere $S^2$ implicitly:
\[ S^2 = \textcolor{eDarkGreen}{F}^{-1}\set{1} = \set{(x,y,z) \in \R^3 : \textcolor{eDarkGreen}{F}(x,y,z) = 1} \]
where the constraint is $\textcolor{eDarkGreen}{F}(x,y,z) = x^2 + y^2 + z^2$.

We determine the normal space $N_{e_3} S^2$ at the North Pole $e_3 = (0,0,1)$:
\[ \nabla \textcolor{eDarkGreen}{F}(0,0,1) = (2x, 2y, 2z) \Big|_{(x,y,z)=(0,0,1)} = (0,0,2) = 2e_3 \]
So, the normal space points straight up along the z-axis:
\[ N_{e_3} S^2 = \Span\set{e_3} = \set{0} \times \set{0} \times \R \]

\begin{center}
\begin{tikzpicture}[scale=2, very thick]
    % Sphere
    \draw[eMediumBlue] (0,0) circle (1cm);
    \draw[eMediumBlue, dashed] (1,0) arc (0:180:1cm and 0.3cm);
    \draw[eMediumBlue] (1,0) arc (0:-180:1cm and 0.3cm);

    % Normal Vector
    \draw[->, eSaddleBrown] (0,0.8) -- (0, 1.8) node[right] {$N_{e_3} S^2$};
    \fill[eSaddleBrown] (0,1) circle (1.5pt) node[below right] {$e_3$};
\end{tikzpicture}
\end{center}
\end{example*}


\end{document}

% \section{Compactness}

% Let $(X, d)$ be a metric space. A subset $K \subseteq X$ is called \textcolor{eDarkRed}{\textbf{compact}} if one of the following equivalent conditions hold:

% \begin{enumerate}
%     \item $K$ is \textcolor{eMediumBlue}{\qt{sequentially compact}}: every sequence $(x_n)_{n \in \N}$ in $K$ has a subsequence that converges to a point in $K$.
%     \item $K$ is \textcolor{eMediumBlue}{\qt{topologically compact}}: every family of open sets $\set{U_i}_{i \in I}$ that cover $K$, has a finite subcover.
% \end{enumerate}

% \begin{center}
% \begin{tikzpicture}[scale=1.2, thick]
%     % The space X
%     \draw[eDarkGray, dashed] (-3, -2) rectangle (5, 3);
%     \node[eDarkGray] at (4.5, 2.5) {$X$};

%     % The compact set K
%     \draw[eDarkRed, fill=eDarkRed!20, rounded corners=15pt] (0,0) -- (1,2) -- (3,1.5) -- (2,-1) -- cycle;
%     \node[eDarkRed] at (1.5, 0.5) {\Large $K$};

%     % An infinite open cover (faded)
%     \foreach \x/\y in {0/0, 0.5/1, 1/1.8, 1.5/0.2, 2/1, 2.5/-0.5, 2.8/1.2, 0.2/1.5, 1.8/-0.8, 0.8/-0.2} {
%         \draw[eDarkGreen!30, fill=eDarkGreen!5, dashed] (\x,\y) circle (0.8cm);
%     }

%     % A finite subcover (highlighted)
%     \draw[eMediumBlue, fill=eMediumBlue!10, fill opacity=0.5, very thick] (0.2, 0.8) circle (1.2cm);
%     \draw[eMediumBlue, fill=eMediumBlue!10, fill opacity=0.5, very thick] (2.2, 0.5) circle (1.3cm);
%     \draw[eMediumBlue, fill=eMediumBlue!10, fill opacity=0.5, very thick] (1.2, -0.2) circle (1cm);
    
%     \node[eMediumBlue] at (-1, 1.5) {$U_{\alpha_1}$};
%     \node[eMediumBlue] at (3.5, 1.2) {$U_{\alpha_2}$};
%     \node[eMediumBlue] at (1.2, -1.4) {$U_{\alpha_3}$};
    
% \end{tikzpicture}
% \end{center}
% \textcolor{eGray}{Visualizing Topological Compactness: Out of an infinite sea of open sets covering $K$, we only actually need a finite number (the blue ones) to completely cover $K$.}

% \vspace{0.5cm}

% \begin{exercises}
% \begin{enumerate}
%     \item Show that every finite subset $\set{x_1, \dots, x_n} \subseteq X$ is compact.
%     \item If $K$ is compact, is it complete?
% \end{enumerate}
% \end{exercises}

% \begin{solution}
% \textbf{1.} Let $\set{U_\alpha}_{\alpha \in I}$ be an open cover of $\set{x_1, \dots, x_n}$. For each $i=1, \dots, n$, we can pick a set $U_{\alpha_i}$ from our cover such that $x_i \in U_{\alpha_i}$. 
% Then the collection $\set{U_{\alpha_1}, \dots, U_{\alpha_n}}$ is a finite subcover. Thus, the set is compact.

% \vspace{0.3cm}
% \textbf{2.} Yes. Let $(x_n)_{n=0}^\infty$ be a Cauchy sequence in $X$. Because $K$ is compact, we can find a convergent subsequence. Since every Cauchy sequence with a convergent subsequence must itself converge to the exact same limit, $(x_n)$ converges in $K$. Thus, $K$ is complete.
% \end{solution}
